% formal/formal.tex
% mainfile: ../perfbook.tex
% SPDX-License-Identifier: CC-BY-SA-3.0

\QuickQuizChapter{chp:Formal Verification}{形式化验证}{qqzformal}
%
\Epigraph{当心上述代码中的 bug；我只是证明了它是正确的，
	  而没有试过。}{Donald Knuth}

\OriginallyPublished{chp:Formal Verification}{Formal Verification}{Linux Weekly News}{PaulEMcKenney2007QRCUspin,PaulEMcKenney2008dynticksRCU,PaulEMcKenney2011ppcmem}

并行算法可能很难编写，而且更难调试。
测试虽然必不可少，但却是不充分的，因为致命的\IXpl{race condition}
可能具有极低的发生概率。
正确性证明可以是有价值的，但归根结底它与原始算法一样
容易出现人为错误。
此外，正确性证明不能指望发现你假设中的错误、需求中的缺陷、
对底层软件或硬件原语的误解，
或者你没有想到要为之构造证明的错误。
这意味着形式化方法永远不能取代测试。
然而，形式化方法可以成为你验证工具箱中有价值的补充。

如果有一个工具能以某种方式定位所有的竞态条件，那将非常有帮助。
存在许多这样的工具，例如，
\cref{sec:formal:State-Space Search}介绍了
通用 state-space search 工具 Promela 和 Spin，
\cref{sec:formal:Special-Purpose State-Space Search}
类似地介绍了专用的 ppcmem 工具，
\cref{sec:formal:Axiomatic Approaches}
考察了一个 axiomatic approach 示例，
\cref{sec:formal:SAT Solvers}
简要概述了 SAT solvers，
\cref{sec:formal:Stateless Model Checkers}
简要概述了 stateless model checkers，
\cref{sec:formal:Summary}
总结了 formal verification 工具在验证并行算法中的使用，
最后
\cref{sec:formal:Choosing a Validation Plan}
讨论了如何决定对给定软件项目应用多少以及什么类型的验证。

% formal/spinhint.html
% mainfile: ../perfbook.tex
% SPDX-License-Identifier: CC-BY-SA-3.0

\section{状态空间搜索}
\label{sec:formal:State-Space Search}
%
\epigraph{探索每条小径 / 走遍你知晓的每条路。}
	 {\emph{Climb Every Mountain}, Rodgers \& Hammerstein}

本节介绍通用的 Promela 和 Spin 工具，
它们可用于对多线程代码进行多种类型的完整状态空间搜索。
对于数据通信协议验证来说，它们也非常有用。
\Cref{sec:formal:Promela and Spin}
介绍了 Promela 和 Spin，包括两个热身练习——
分别验证非原子递增和原子递增。
\Cref{sec:formal:How to Use Promela}
描述了 Promela 的使用方法，包括示例命令行以及
Promela 语法与 C 语法的比较。
\Cref{sec:formal:Promela Example: Locking}
展示了如何使用 Promela 验证锁机制，
\cref{sec:formal:Promela Example: QRCU}
使用 Promela 验证一种名为``QRCU''的特殊 RCU 实现，
最后
\cref{sec:formal:Promela Parable: dynticks and Preemptible RCU}
将 Promela 应用于 RCU dyntick-idle 实现的早期版本。

\subsection{Promela 和 Spin}
\label{sec:formal:Promela and Spin}

\IX{Promela} 是一种旨在帮助验证协议的语言，但也可以用于验证小型并行算法。
你需要用类似于 C 语言的 Promela 来重新编写你的算法和正确性约束，
然后使用 \IX{Spin} 将其翻译为可以编译和运行的 C 程序。
生成的程序会对你的算法执行完整的状态空间搜索，
要么确认通过你在 Promela 程序中关联的断言，要么找到反例。

这种完整的状态搜索非常强大，但也是一把双刃剑。
如果算法太复杂，或者 Promela 实现得不够仔细，
那么状态空间可能大到内存无法容纳。
此外，即使有充足的内存，状态空间搜索的运行时间
也可能超过宇宙的预期寿命。
因此，请在复杂但紧凑的并行算法中使用此工具。
轻率地试图将其应用于中等规模的算法（更不用说整个 Linux 内核）
将会以失败告终。

Promela 和 Spin 可以从
\url{https://spinroot.com/spin/whatispin.html} 下载。

上述网站还提供了 Gerard Holzmann 关于 Promela 和 Spin 的优秀
著作~\cite{Holzmann03a}的链接，
可以通过以下地址进行在线搜索：
\url{https://www.spinroot.com/spin/Man/index.html}。

本节其余部分描述如何使用 Promela 调试并行算法，
从简单示例开始，然后逐步过渡到更复杂的使用场景。

\subsubsection{热身：
			非原子递增}
\label{sec:formal:Warm-Up: Non-Atomic Increment}

\begin{fcvref}[ln:formal:promela:increment:whole]
\Cref{lst:formal:Promela Code for Non-Atomic Increment}
演示了教科书中所讲的非原子递增导致的\IX{竞态条件}。
\Clnref{nprocs} 定义了要运行的进程数量（我们将改变此值
以观察对状态空间的影响），\clnref{count} 定义了计数器，
\clnref{prog} 用于实现
\clnrefrange{assert:b}{assert:e} 上的断言。

\begin{listing}
\input{CodeSamples/formal/promela/increment@whole.fcv}
\caption{非原子递增的 Promela 代码}
\label{lst:formal:Promela Code for Non-Atomic Increment}
\end{listing}

\Clnrefrange{proc:b}{proc:e} 定义了一个非原子地递增计数器的过程。
参数 \co{me} 是进程编号，由后面的初始化代码进行设置。
由于简单的 Promela 语句各自被假定为原子的，
我们必须将递增操作拆分为
\clnrefrange{incr:b}{incr:e} 上的两条语句。
\clnref{setprog} 上的赋值标记该过程的完成。
由于 Spin 系统会搜索整个状态空间，包括所有可能的状态序列，
因此不需要传统压力测试中使用的循环。

\Clnrefrange{init:b}{init:e} 是初始化代码块，它最先被执行。
\Clnrefrange{doinit:b}{doinit:e} 实际进行初始化操作，
而 \clnrefrange{assert:b}{assert:e} 执行断言。
两者都是原子块，以避免不必要地增加状态空间：
因为它们不属于算法本身的一部分，
将它们标记为原子的不会导致验证覆盖率的缺失。

\clnrefrange{dood1:b}{dood1:e} 上的 \co{do-od} 构造实现了一个 Promela 循环，
可以将其理解为一个 C 语言的 \co{for (;;)} 循环，
其中包含一个允许在 case 标签中使用表达式的 \co{switch} 语句。
条件块（以 \co{::} 为前缀）会被非确定性地扫描，
不过在本例中，在任何给定时刻只有一个条件可能成立。
\co{do-od} 的第一个块从
\clnrefrange{block1:b}{block1:e}，
初始化第 i 个递增器的进度单元，运行第 i 个递增器的进程，
然后递增变量 \co{i}。
\co{do-od} 的第二个块在 \clnref{block2} 上，
在这些进程启动后退出循环。

\clnrefrange{assert:b}{assert:e} 上的原子块也包含
一个类似的 \co{do-od} 循环，用于累加进度计数器。
\clnref{assert} 上的 \co{assert()} 语句验证：
如果所有进程都已完成，则所有计数都已正确记录。
\end{fcvref}

你可以按如下方式构建并运行此程序：

\begin{VerbatimU}
spin -a increment.spin      # Translate the model to C
cc -DSAFETY -o pan pan.c    # Compile the model
./pan                       # Run the model
\end{VerbatimU}

\begin{listing}
\VerbatimInput[numbers=none,fontsize=\scriptsize]{CodeSamples/formal/promela/increment.spin.lst}
\vspace*{-9pt}
\caption{非原子递增的 Spin 输出}
\label{lst:formal:Non-Atomic Increment Spin Output}
\end{listing}

这将产生如
\cref{lst:formal:Non-Atomic Increment Spin Output}
所示的输出。
第一行告诉我们断言被违反了（考虑到非原子递增，这在预期之中！）。
第二行中的跟踪（trail）文件描述了断言是如何被违反的。
``Warning''行重申在我们的模型中，并非所有事情都正确完成。
第二段描述了执行的状态搜索类型，
本例中是检查断言违反和无效终止状态。
第三段给出了状态大小统计信息：
这个小模型只有 45 个状态。
最后一行显示了内存使用情况。

跟踪（trail）文件可以被整理为如下所示的人类可读格式：

\begin{VerbatimU}
spin -t -p increment.spin
\end{VerbatimU}

\begin{listing*}
\ebresizeverb{.9}{
\VerbatimInput[numbers=none,fontsize=\scriptsize]{CodeSamples/formal/promela/increment.spin.trail.lst}}
\IfEbookSize{\vspace*{7pt}}{\vspace*{-9pt}}
\caption{非原子递增错误轨迹}
\label{lst:formal:Non-Atomic Increment Error Trail}
\end{listing*}

它给出了如
\cref{lst:formal:Non-Atomic Increment Error Trail}
所示的输出。
正如我们所见，初始化块的第一部分创建了两个递增器进程，
其中第一个进程获得了计数器的值，然后两个进程都递增并存储了它，
从而丢失了一次计数。
随后断言被触发，之后显示了全局状态。

\subsubsection{热身：
			原子递增}
\label{sec:formal:Warm-Up: Atomic Increment}

像\cref{lst:formal:Promela Code for Atomic Increment}那样，
通过将递增过程的主体放入原子块中，可以简单地修复此示例。
由于 Promela 语句是原子的，一个简单的方法是用
\co{counter = counter + 1} 语句来替换递增语句对。
无论哪种方式，运行修改后的模型都会给出无错误的状态空间遍历，
如\cref{lst:formal:Atomic Increment Spin Output}所示。

\begin{listing}
\input{CodeSamples/formal/promela/atomicincrement@incrementer.fcv}
\caption{原子递增的 Promela 代码}
\label{lst:formal:Promela Code for Atomic Increment}
\end{listing}

\begin{listing}
\VerbatimInput[numbers=none,fontsize=\scriptsize]{CodeSamples/formal/promela/atomicincrement.spin.lst}
\vspace*{-9pt}
\caption{原子递增的 Spin 输出}
\label{lst:formal:Atomic Increment Spin Output}
\end{listing}

\Cref{tab:advsync:Memory Usage of Increment Model}
显示了作为所建模的递增器数量（通过重新定义 \co{NUMPROCS}）
的函数的状态数和内存消耗：

\begin{table}
\rowcolors{1}{}{lightgray}
\small
\renewcommand*{\arraystretch}{1.2}
\centering
\begin{tabular}{S[table-format = 1.0]S[table-format = 7.0]S[table-format = 3.1]}
	\toprule
	\multicolumn{1}{l}{递增器数量} &
		\multicolumn{1}{r}{状态数} &
			\multicolumn{1}{r}{总内存使用量 (MB)} \\
	\midrule
	1 &		        11 &        128.7 \\
	2 &		        52 &        128.7 \\
	3 &		       372 &        128.7 \\
	4 &		     3 496 &        128.9 \\
	5 &		    40 221 &        131.7 \\
	6 &		   545 720 &        174.0 \\
	7 &		 8 521 446 &        881.9 \\
	\bottomrule
\end{tabular}
\caption{递增模型的内存使用量}
\label{tab:advsync:Memory Usage of Increment Model}
\end{table}

因此，不必要地运行大型模型的前景看起来不容乐观，尽管
882\,MB 完全在现代台式机和笔记本电脑的限制范围内。

有了这个示例的经验，让我们更仔细地看看
用于分析 Promela 模型的命令，然后看看更精细的示例。

\subsection{如何使用 Promela}
\label{sec:formal:How to Use Promela}

给定一个源文件 \path{qrcu.spin}，可以使用以下命令：

\begin{description}[style=nextline]
\item	[\tco{spin -a qrcu.spin}]
	创建一个完整搜索状态机的文件 \path{pan.c}。
\item	[\tco{cc -DSAFETY [-DCOLLAPSE] [-DMA=N] -o pan pan.c}]
	编译生成的状态机搜索程序。
	\co{-DSAFETY} 生成适用于仅包含断言（以及可能的 \co{never} 语句）
	情况下的优化。
	如果你有活性、公平性或前向进展检查，
	可能需要在不使用 \co{-DSAFETY} 的情况下编译。
	如果你在可以使用 \co{-DSAFETY} 时没有使用它，
	程序会提示你。

	\co{-DSAFETY} 产生的优化大大加快了速度，
	所以应尽可能使用它。
	不能使用 \co{-DSAFETY} 的一个示例场景是
	通过 \co{-DNP} 检查\IXpl{livelock}（又称``非进展循环''）时。

	可选的 \co{-DCOLLAPSE} 标志生成状态向量压缩模式的代码。

	另一个可选标志 \co{-DMA=N} 生成一种较慢但更激进的
	状态空间内存压缩模式的代码。
\item	[\tco{./pan [-mN] [-wN]}]
	这将实际搜索状态空间。
	即使对于非常小的状态机，状态数也可以达到数千万，
	因此你需要一台具有大内存的机器。
	例如，具有 3 个更新者和 2 个读者的 \path{qrcu.spin}
	即使使用 \co{-DCOLLAPSE} 标志也需要 10.5\,GB 内存。

	如果你看到 \co{./pan} 显示的消息：
	``\co{error: max search depth too small}''，你需要通过
	\co{-mN} 选项增加最大深度以进行完整搜索。
	默认值为 \co{-m10000}。

	\co{-wN} 选项指定哈希表（hash table）大小。
	完整状态空间搜索的默认值为 \co{-w24}。\footnote{
		截至 Spin 版本 6.4.6 和 6.4.8。
		在 2011 年 7 月 10 日的 Spin 在线手册中，
		穷举搜索模式的默认值据说是 \co{-w19}，
		这与实际行为不符。}

	如果你不确定你的机器是否有足够的内存，
	在一个窗口中运行 \co{top}，在另一个窗口中运行 \co{./pan}。
	将焦点保持在 \co{./pan} 窗口上，以便在需要时快速终止执行。
	一旦 CPU 时间明显低于 100\pct，就终止 \co{./pan}。
	如果你已将焦点从运行 \co{./pan} 的窗口移开，
	可能需要等待很长时间，窗口系统才能获得足够的内存来为你做任何事情。

	另一种避免内存耗尽的选择是
	\co{-DMEMLIM=N} 编译器标志。
	\co{-DMEMLIM=2000} 将设置 2\,GB 的最大值。

	不要忘记捕获输出，特别是在远程机器上工作时。

	如果你的模型包含前向进展检查，你可能需要
	通过 \co{./pan} 的 \co{-f} 命令行参数启用``弱公平性''。
	如果你的前向进展检查涉及 \co{accept} 标签，
	你还需要 \co{-a} 参数。
	% forward reference to model: formal.2009.02.19a in
	% /home/linux/git/userspace-rcu/formal-model.
\item	[\tco{spin -t -p qrcu.spin}]
	给定遇到错误的运行所输出的 \co{trail} 文件，
	输出导致该错误的步骤序列。
	\co{-g} 标志还会包含已更改的全局变量的值，
	\co{-l} 标志还会包含已更改的局部变量的值。
\end{description}

\subsubsection{Promela 的特殊之处}
\label{sec:formal:Promela Peculiarities}

尽管所有计算机语言都有潜在的相似性，
但 Promela 仍然会给习惯于用 C、C++ 或 Java 编程的人带来一些意外。

\begin{enumerate}
\item	在 C 中，\qco{;} 终止一条语句。
	在 Promela 中，它将语句进行分隔。
	幸运的是，较新版本的 Spin 已经对此比较宽容了。
\item	Promela 的循环构造 \co{do} 语句需要循环条件。
	该 \co{do} 语句类似于 if-then-else 语句的循环版本。
\item	在 C 的 \co{switch} 语句中，如果没有匹配的情况，
	整个语句被忽略。
	在 Promela 中，其等价物令人困惑地被称为 \co{if}，
	如果没有相应的守卫表达式匹配，将得到一个错误，
	而且没有相应的可识别的错误消息。
	因此，如果错误输出指向无辜的代码行，
	请检查你是否在 \co{if} 或 \co{do} 语句中遗漏了某个条件。
\item	在 C 中创建压力测试时，通常会将可疑操作相互放在一起运行。
	在 Promela 中，取而代之的是只运行一次操作，因为 Promela
	会搜索出该操作所有可能的结果。
	有时你确实需要在 Promela 中循环操作，例如，
	当多个操作交叉时，但这样做会大大增加状态空间的大小。
\item	在 C 中，最简单的做法是维护一个循环计数器来跟踪进度和终止循环。
	在 Promela 中，循环计数器必须像躲避瘟疫一样避免使用，
	因为它们会导致状态空间迅速膨胀。
	但从另一方面来说，在 Promela 中使用无限循环是没有问题的，
	只要循环里面没有变量单调递增或递减——Promela 会自行计算出
	循环到底需要多少次遍历才有意义，
	并自动删除超出该点的不必要执行。
\item	在 C 压力测试代码中，使用每任务控制变量通常是明智的。
	它们读取开销低，对调试测试代码也很有帮助。
	在 Promela 中，仅仅在没有其他替代方案时，才使用每任务变量。
	要明白这一点，考虑一个 5 任务的验证，
	每一位表示一个任务的完成。
	这会导致 32 个状态。
	相比之下，一个简单的计数器只有六个状态，
	状态数整整少了五倍多。
	五倍看起来可能不是什么问题，但如果你在为一个拥有超过
	1.5 亿个状态、消耗超过 10\,GB 内存的验证程序而挣扎时，
	就不会这样想了！
\item	无论是在 C 压力测试代码还是 Promela 中，最具挑战性的事情之一
	是编写好的断言。
	Promela 还允许 \co{never} 声明，它的作用类似于
	在每一行代码之间复制的断言。
\item	在 Promela 中，分而治之对于将状态空间保持在可控范围内非常有用。
	将一个大型模型分成大致相等的两半，
	会将状态空间减少到大约整体的平方根。
	例如，一个百万状态的组合模型可能缩减为
	一对千状态的模型。
	Promela 不仅能以更少的内存更快地处理两个较小的模型，
	而且两个较小的算法也更容易让人理解。
\end{enumerate}


\subsubsection{Promela 编码技巧}
\label{sec:formal:Promela Coding Tricks}

Promela 被设计用于分析协议，因此将其用于并行程序算是有点滥用。
以下技巧可以帮助你安全地用好 Promela：

\begin{enumerate}
\item	内存乱序。
	假设你有一对语句将全局变量 \co{x} 和 \co{y}
	复制到局部变量 \co{r1} 和 \co{r2}，其中顺序很重要（例如，没有锁保护），
	但你没有\IXpl{memory barrier}。
	这可以在 Promela 中用如下方法建模：

\begin{VerbatimN}[samepage=true]
if
:: 1 -> r1 = x;
        r2 = y
:: 1 -> r2 = y;
        r1 = x
fi
\end{VerbatimN}

	\co{if} 语句的两个分支会被非确定性地选取，因为它们都是可用的。
	由于进行了完整的状态空间搜索，在所有情况下
	\emph{两种}选择最终都会被覆盖。

	当然，过度使用这个技巧会导致你的状态空间迅速增大。
	此外，它要求你预料到可能的乱序情况。

\item	状态压缩。
	如果你有复杂的断言，请在 \co{atomic} 块内求值。
	毕竟，它们不是算法的一部分。
	一个复杂断言的例子（后面将更详细地讨论）
	如\cref{lst:formal:Complex Promela Assertion}所示。

	没有理由非原子地执行这个断言，
	因为它实际上不是算法的一部分。
	因为每条语句都会增加状态数量，我们可以通过将其
	放入 \co{atomic} 块中来减少无用状态的数量，
	如\cref{lst:formal:Atomic Block for Complex Promela Assertion}所示。

\item	Promela 不提供函数。
	必须使用 C 预处理器宏来代替。
	但是，要小心使用它们以避免状态组合快速增加。
\end{enumerate}

\begin{listing}
\begin{VerbatimL}
i = 0;
sum = 0;
do
:: i < N_QRCU_READERS ->
	sum = sum + (readerstart[i] == 1 &&
	             readerprogress[i] == 1);
	i++
:: i >= N_QRCU_READERS ->
	assert(sum == 0);
	break
od
\end{VerbatimL}
\caption{复杂的 Promela 断言}
\label{lst:formal:Complex Promela Assertion}
\end{listing}

\begin{listing}
\begin{VerbatimL}
atomic {
	i = 0;
	sum = 0;
	do
	:: i < N_QRCU_READERS ->
		sum = sum + (readerstart[i] == 1 &&
		             readerprogress[i] == 1);
		i++
	:: i >= N_QRCU_READERS ->
		assert(sum == 0);
		break
	od
}
\end{VerbatimL}
\caption{复杂 Promela 断言的原子块}
\label{lst:formal:Atomic Block for Complex Promela Assertion}
\end{listing}

现在我们准备好见识更复杂的示例了。

\subsection{Promela 示例：
			     锁}
\label{sec:formal:Promela Example: Locking}

\begin{fcvref}[ln:formal:promela:lock:whole]
由于锁通常是有用的，\path{lock.h} 中提供了 \co{spin_lock()} 和
\co{spin_unlock()} 宏，可以从多个 Promela 模型中包含，
如\cref{lst:formal:Promela Code for Spinlock}所示。
\co{spin_lock()} 宏包含一个无限 \co{do-od} 循环，
跨越 \clnrefrange{dood:b}{dood:e}，
由 \clnref{one} 上的单一守卫表达式``1''驱动。
该循环的主体是一个包含 \co{if-fi} 语句的原子块。
\co{if-fi} 构造与 \co{do-od} 构造类似，
只是它只执行一次而不是循环。
如果锁在 \clnref{notheld} 上未被持有，则
\clnref{acq} 获取它，
\clnref{break} 跳出外层 \co{do-od} 循环（同时也退出原子块）。
另一方面，如果锁在 \clnref{held} 上已被持有，
我们什么都不做（\co{skip}），并跳出 \co{if-fi} 和原子块，
以便再次通过外层循环，重复直到锁可用。
\end{fcvref}

\begin{listing}
\input{CodeSamples/formal/promela/lock@whole.fcv}
\caption{自旋锁（spinlock）的 Promela 代码}
\label{lst:formal:Promela Code for Spinlock}
\end{listing}

\co{spin_unlock()} 宏简单地将锁标记为不再被持有。

请注意，不需要内存屏障（memory barrier），因为 Promela 假设强序。
在任何给定的 Promela 状态中，所有处理器对当前状态以及
导致我们到达当前状态的状态变化顺序都达成一致。
这类似于少数计算机系统（如 1990 年代的 MIPS 和 PA-RISC）
所使用的``\IXalth{sequentially consistent}{sequential}{memory consistency}''
\IX{memory model}。
如前面的提示所述，并且将在后面的示例中看到，
弱内存排序必须被显式编码。

\begin{listing}
\input{CodeSamples/formal/promela/lock@spin.fcv}
\caption{测试自旋锁的 Promela 代码}
\label{lst:formal:Promela Code to Test Spinlocks}
\end{listing}

\begin{fcvref}[ln:formal:promela:lock:spin]
这些宏通过\cref{lst:formal:Promela Code to Test Spinlocks}
中所示的 Promela 代码进行测试。
此代码与用于测试递增的代码类似，
锁定进程的数量由 \clnref{nlockers} 上的
\co{N_LOCKERS} 宏定义来定义。
互斥锁（exclusive lock）本身在 \clnref{mutex} 上定义，
用于跟踪锁持有者的数组在 \clnref{array} 上，
\clnref{sum} 由断言代码使用以验证只有一个进程持有锁。
\end{fcvref}

\begin{fcvref}[ln:formal:promela:lock:spin:locker]
锁定器进程在 \clnrefrange{b}{e} 上，简单地无限循环：
在 \clnref{lock} 上获取锁，在 \clnref{claim} 上声明持有，
在 \clnref{unclaim} 上取消声明，在 \clnref{unlock} 上释放锁。
\end{fcvref}

\begin{fcvref}[ln:formal:promela:lock:spin:init]
\clnrefrange{b}{e} 上的 init 块在 \clnref{array} 上初始化当前锁定器的
\co{havelock} 数组项，在 \clnref{start} 上启动当前锁定器，
并在 \clnref{next} 上前进到下一个锁定器。
一旦所有锁定器进程被创建，\co{do-od} 循环
移动到 \clnref{chkassert}，用于检查断言。
\Clnref{sum,j} 初始化控制变量，
\clnrefrange{atm:b}{atm:e} 原子地对 \co{havelock} 数组项求和，
\clnref{assert} 是断言，\clnref{break} 退出循环。
\end{fcvref}

我们可以通过将
\cref{lst:formal:Promela Code for Spinlock,%
lst:formal:Promela Code to Test Spinlocks} 中的两个代码片段
分别放入名为 \path{lock.h} 和 \path{lock.spin} 的文件中，
然后运行以下命令来运行此模型：

\begin{VerbatimU}
spin -a lock.spin
cc -DSAFETY -o pan pan.c
./pan
\end{VerbatimU}

\begin{listing}
\VerbatimInput[numbers=none,fontsize=\scriptsize]{CodeSamples/formal/promela/lock.spin.lst}
\vspace*{-9pt}
\caption{自旋锁测试的输出}
\label{lst:formal:Output for Spinlock Test}
\end{listing}

输出将看起来类似于
\cref{lst:formal:Output for Spinlock Test}所示。
如预期的那样，此运行没有断言失败（\qco{errors: 0}）。

\QuickQuizSeries{%
\QuickQuizB{
	为什么 locker 中有一条未到达的语句？
	毕竟，这不是\emph{完整的}状态空间搜索吗？
}\QuickQuizAnswerB{
	locker 进程是一个无限循环，因此控制永远不会
	到达该进程的末尾。
	但是，由于没有单调递增的变量，
	Promela 能够用少量状态来建模这个无限循环。
}\QuickQuizEndB
%
\QuickQuizE{
	这个示例有哪些 Promela 代码风格问题？
}\QuickQuizAnswerE{
	有几个：
	\begin{enumerate}
	\item	\co{sum} 的声明应该移到 init 块内部，
		因为它在其他地方没有被使用。
	\item	断言代码应该移到初始化循环之外。
		然后可以将初始化循环放在原子块中，
		大大减少状态空间（减少多少？）。
	\item	覆盖断言代码的原子块应该扩展到包含
		\co{sum} 和 \co{j} 的初始化，
		并且也覆盖断言。
		这也减少了状态空间（又减少了多少？）。
	\end{enumerate}
}\QuickQuizEndE
}

\subsection{Promela 示例：
			     QRCU}
\label{sec:formal:Promela Example: QRCU}

这个最终示例演示了在 Oleg Nesterov 的
QRCU~\cite{OlegNesterov2006QRCU,OlegNesterov2006aQRCU}
上实际使用 Promela 的情况，
但进行了修改以加速 \co{synchronize_qrcu()} 的快速路径。

但首先需要搞清楚，QRCU 是什么？

QRCU 是 SRCU~\cite{PaulEMcKenney2006c} 的一个变体，
它以较高的读开销（对全局变量的原子递增和递减）换取
极低的宽限期延迟。
如果没有读者，宽限期将在不到一微秒内被检测到，
而大多数其他 RCU 实现的宽限期延迟则为数毫秒。

\begin{enumerate}
\item	有一个 \co{qrcu_struct} 定义了一个 QRCU 域。
	与 SRCU 类似（不同于其他 RCU 变体），QRCU 的作用
	不是全局的，而是集中在指定的 \co{qrcu_struct} 上。
\item	有 \co{qrcu_read_lock()} 和 \co{qrcu_read_unlock()}
	原语用于界定 QRCU 读端临界区（critical section）。
	必须将相应的 \co{qrcu_struct} 传递给这些原语，
	并且 \co{qrcu_read_lock()} 的返回值
	必须传递给 \co{qrcu_read_unlock()}。

	例如：

\begin{VerbatimU}
idx = qrcu_read_lock(&my_qrcu_struct);
/* read-side critical section. */
qrcu_read_unlock(&my_qrcu_struct, idx);
\end{VerbatimU}

\item	有一个 \co{synchronize_qrcu()} 原语，它会阻塞直到
	所有先前存在的 QRCU 读端临界区完成，
	但与 SRCU 的 \co{synchronize_srcu()} 类似，QRCU 的
	\co{synchronize_qrcu()} 只需等待使用相同 \co{qrcu_struct}
	的那些读端临界区。

	例如，\co{synchronize_qrcu(&your_qrcu_struct)}
	\emph{不}需要等待前面的 QRCU 读端临界区。
	相反，\co{synchronize_qrcu(&my_qrcu_struct)}
	\emph{需要}等待，因为它共享相同的 \co{qrcu_struct}。
\end{enumerate}

已经产生了一个 QRCU 的 Linux 内核补丁~\cite{PaulMcKenney2007QRCUpatch}，
但它不太可能被包含在 Linux 内核中。

\begin{listing}
\input{CodeSamples/formal/promela/qrcu@gvar.fcv}
\caption{QRCU 全局变量}
\label{lst:formal:QRCU Global Variables}
\end{listing}

回到 QRCU 的 Promela 代码，全局变量如
\cref{lst:formal:QRCU Global Variables}所示。
这个示例使用了锁并包含了 \path{lock.h}。
读者和写者的数量都可以通过两个 \co{#define} 语句来改变，
这给了我们不止一种而是两种造成组合爆炸的方式。
\co{idx} 变量控制 \co{ctr} 数组的两个元素中哪一个将被读者使用，
\co{readerprogress} 变量允许断言判断所有读者何时完成
（因为直到每个已存在的读者完成其 QRCU 读端临界区之前，
QRCU 更新都不允许完成）。
\co{readerprogress} 数组元素的值如下，
表示相应读者的状态：

\begin{enumerate}[label={\arabic*}:,start=0,itemsep=0pt]
\item	尚未开始。
\item	在 QRCU 读端临界区内。
\item	已完成 QRCU 读端临界区。
\end{enumerate}

最后，\co{mutex} 变量用于序列化更新者的慢速路径。

\begin{listing}
\input{CodeSamples/formal/promela/qrcu@reader.fcv}
\caption{QRCU 读者进程}
\label{lst:formal:QRCU Reader Process}
\end{listing}

\begin{fcvref}[ln:formal:promela:qrcu:reader]
QRCU 读者由\cref{lst:formal:QRCU Reader Process}中所示的
\co{qrcu_reader()} 进程建模。
一个 \co{do-od} 循环跨越 \clnrefrange{do}{od}，
在 \clnref{one} 上有一个值为``1''的单一守卫，
使其成为无限循环。
\Clnref{curidx} 捕获全局索引的当前值，
\clnrefrange{atm:b}{atm:e} 在其值非零时原子地递增它
（并跳出无限循环）（\co{atomic_inc_not_zero()}）。
\Clnref{cs:entry} 标记进入 RCU 读端临界区，
\clnref{cs:exit} 标记退出该临界区，
这两行都是为了后面将遇到的 \co{assert()} 语句服务的。
\Clnref{atm:dec} 原子地递减我们之前递增的同一计数器，
从而退出 RCU 读端临界区。
\end{fcvref}

\begin{listing}
\input{CodeSamples/formal/promela/qrcu@sum_unordered.fcv}
\caption{QRCU 无序求和}
\label{lst:formal:QRCU Unordered Summation}
\end{listing}

\begin{fcvref}[ln:formal:promela:qrcu:sum_unordered]
\cref{lst:formal:QRCU Unordered Summation}
中所示的 C 预处理器宏对一对计数器求和，
以模拟弱内存排序。
\Clnrefrange{fetch:b}{fetch:e} 获取其中一个计数器，
\clnref{sum_other} 获取另一个并求和。
原子块由单个 \co{do-od} 语句组成。
这个 \co{do-od} 语句（跨越 \clnrefrange{do}{od}）不寻常之处在于
它包含两个在 \clnref{g1,g2} 上带有守卫的无条件分支，
这导致 Promela 非确定性地选择其中一个（但同样，
完整的状态空间搜索使 Promela 最终在每个适用情况下做出所有可能的选择）。
第一个分支获取第零个计数器并将 \co{i} 设置为 1（以便
\clnref{sum_other} 获取第一个计数器），而第二个分支相反，
获取第一个计数器并将 \co{i} 设置为 0（以便
\clnref{sum_other} 获取第二个计数器）。
\end{fcvref}

\QuickQuiz{
	有没有更直接的方式来编写 \co{do-od} 语句？
}\QuickQuizAnswer{
	有。
	将其替换为 \co{if-fi} 并删除两个 \co{break} 语句。
}\QuickQuizEnd

\begin{listing}
\ebresizeverb{.88}{\input{CodeSamples/formal/promela/qrcu@updater.fcv}}
\caption{QRCU 更新者进程}
\label{lst:formal:QRCU Updater Process}
\end{listing}

\begin{fcvref}[ln:formal:promela:qrcu:updater]
有了 \co{sum_unordered} 宏，我们现在可以继续看
\cref{lst:formal:QRCU Updater Process}中所示的更新端进程。
更新端进程无限重复，相应的 \co{do-od} 循环
跨越 \clnrefrange{do}{od}。
每次循环首先在 \clnrefrange{atm1:b}{atm1:e} 上
将全局 \co{readerprogress} 数组快照到局部 \co{readerstart} 数组中。
此快照将用于 \clnref{assert} 上的断言。
\Clnref{sum_unord} 调用 \co{sum_unordered}，然后
\clnrefrange{reinvoke:b}{reinvoke:e}
在快速路径可能可用时重新调用 \co{sum_unordered}。

\Clnrefrange{slow:b}{slow:e} 在需要时执行慢速路径代码，
\clnref{acq,rel} 获取和释放更新端锁，
\clnrefrange{flip_idx:b}{flip_idx:e} 翻转索引，
\clnrefrange{wait:b}{wait:e} 等待所有先前存在的读者完成。

\Clnrefrange{atm2:b}{atm2:e} 然后将 \co{readerprogress}
数组中的当前值与 \co{readerstart} 数组中收集的值进行比较，
如果在此更新之前开始的任何读者仍在进行中，
则强制断言失败。
\end{fcvref}

\QuickQuizSeries{%
\QuickQuizB{
	\begin{fcvref}[ln:formal:promela:qrcu:updater]
	为什么在 \clnrefrange{atm1:b}{atm1:e}
	和 \clnrefrange{atm2:b}{atm2:e} 上有原子块，
	而这些原子块中的操作在当前任何生产级微处理器上
	都没有原子实现？
	\end{fcvref}
}\QuickQuizAnswerB{
	因为这些操作仅用于断言。
	它们不是算法本身的一部分。
	因此将它们标记为原子的没有害处，
	这样标记大大减少了 Promela 模型必须搜索的状态空间。
}\QuickQuizEndB
%
\QuickQuizE{
	\begin{fcvref}[ln:formal:promela:qrcu:updater]
	\clnrefrange{reinvoke:b}{reinvoke:e}
	上的计数器重新求和\emph{真的}必要吗？
	\end{fcvref}
}\QuickQuizAnswerE{
	是的。
	要理解这一点，请删除这些行并运行模型。

	或者，考虑以下步骤序列：

	\begin{enumerate}
	\item	一个进程在其 RCU 读端临界区内，
		因此 \co{ctr[0]} 的值为零，
		\co{ctr[1]} 的值为二。
	\item	一个更新者开始执行，看到计数器之和为二，
		因此快速路径不能执行。
		它因此获取了锁。
	\item	第二个更新者开始执行，获取了 \co{ctr[0]} 的值，
		为零。
	\item	第一个更新者将 \co{ctr[0]} 加一，翻转索引
		（现在变为零），然后从 \co{ctr[1]} 减一
		（现在变为一）。
	\item	第二个更新者获取 \co{ctr[1]} 的值，
		现在为一。
	\item	第二个更新者现在错误地得出结论认为
		可以安全地走快速路径，尽管
		原始读者尚未完成。
	\end{enumerate}
}\QuickQuizEndE
}

\begin{listing}
\input{CodeSamples/formal/promela/qrcu@init.fcv}
\caption{QRCU 初始化进程}
\label{lst:formal:QRCU Initialization Process}
\end{listing}

\begin{fcvref}[ln:formal:promela:qrcu:init]
剩下的就是\cref{lst:formal:QRCU Initialization Process}中所示的
初始化块。
这个块只是在 \clnrefrange{i_ctr:b}{i_ctr:e} 上初始化计数器对，
在 \clnrefrange{spn_r:b}{spn_r:e} 上生成读者进程，
在 \clnrefrange{spn_u:b}{spn_u:e} 上生成更新者进程。
这些都在一个原子块内完成以减少状态空间。
\end{fcvref}

\subsubsection{运行 QRCU 示例}
\label{sec:formal:Running the QRCU Example}

要运行 QRCU 示例，请将上一节中的代码片段合并到
一个名为 \path{qrcu.spin} 的文件中，并将
\co{spin_lock()} 和 \co{spin_unlock()} 的定义放入
名为 \path{lock.h} 的文件中。
然后使用以下命令构建并运行 QRCU 模型：

\begin{VerbatimU}
spin -a qrcu.spin
cc -DSAFETY [-DCOLLAPSE] -o pan pan.c
./pan [-mN]
\end{VerbatimU}

\begin{table}
\centering
\begin{threeparttable}
\rowcolors{1}{}{lightgray}
\renewcommand*{\arraystretch}{1.2}
\footnotesize
\begin{tabular}{S[table-format = 1.0]S[table-format = 1.0]S[table-format = 9.0]
		S[table-format = 6.0]S[table-format = 5.1]}
	\toprule
	\multicolumn{1}{r}{更新者} &
	    \multicolumn{1}{r}{读者} &
		\multicolumn{1}{r}{状态数} &
		    \multicolumn{1}{r}{深度} &
			\multicolumn{1}{r}{内存 (MB)\tnote{a}} \\
	\midrule
	1 & 1 &         376 &      95 &    128.7 \\
	1 & 2 &       6 177 &     218 &    128.9 \\
	1 & 3 &      99 728 &     385 &    132.6 \\
	2 & 1 &      29 399 &     859 &    129.8 \\
	2 & 2 &   1 071 181 &   2 352 &    169.6 \\
	2 & 3 &  33 866 736 &  12 857 &  1 540.8 \\
	3 & 1 &   2 749 453 &  53 809 &    236.6 \\
	3 & 2 & 186 202 860 & 328 014 & 10 483.7 \\
	\bottomrule
\end{tabular}
\begin{tablenotes}
	\item [a] 使用编译器标志 \co{-DCOLLAPSE} 获得。
\end{tablenotes}
\end{threeparttable}
\caption{QRCU 模型的内存使用量}
\label{tab:advsync:Memory Usage of QRCU Model}
\end{table}

输出显示此模型通过了
\cref{tab:advsync:Memory Usage of QRCU Model}中所示的所有案例。
如果能运行三个读者和三个更新者就好了，
然而，简单外推表明这将需要大约半太字节的内存。
怎么办？

事实证明，当 \co{./pan} 耗尽内存时会给出建议，
例如，当试图运行三个读者和三个更新者时：

\begin{VerbatimU}
hint: to reduce memory, recompile with
  -DCOLLAPSE # good, fast compression, or
  -DMA=96   # better/slower compression, or
  -DHC # hash-compaction, approximation
  -DBITSTATE # supertrace, approximation
\end{VerbatimU}

让我们尝试建议的编译器标志 \co{-DMA=N}，
它生成以大大增加搜索开销为代价来激进压缩状态空间的代码。
所需命令如下：

\begin{VerbatimU}
spin -a qrcu.spin
cc -DSAFETY -DMA=96 -O2 -o pan pan.c
./pan -m20000000
\end{VerbatimU}

这里，深度限制 20,000,000 比从简单外推推导出的预期深度
大一个数量级。
虽然这增加了前期内存使用，但避免了由于过紧的
深度限制导致的不完整搜索而浪费漫长的运行。
此运行在一台 \Power{9} 服务器上花费了略多于 3 天。
结果如
\cref{lst:formal:spinhint:3 Readers 3 Updaters QRCU Spin Output with -DMA=96}所示。
此 Spin 运行成功完成，总内存使用量仅为 6.5\,GB，
比 \co{-DCOLLAPSE} 使用的约半太字节低了近两个数量级。

\begin{listing}
\VerbatimInput[numbers=none,fontsize=\scriptsize]{CodeSamples/formal/promela/qrcu.spin.33ma.lst}
\vspace*{-9pt}
\caption{使用 \co{-DMA=96} 的 3 读者 3 更新者 QRCU Spin 输出}
\label{lst:formal:spinhint:3 Readers 3 Updaters QRCU Spin Output with -DMA=96}
\end{listing}

\QuickQuiz{
	0.48\pct\ 的压缩率对应于状态占用内存
	200 比 1 的减少！
	状态空间搜索\emph{真的}是穷举的吗？？？
}\QuickQuizAnswer{
	根据 Spin 的文档，是的，它是穷举的。

\begin{listing}
\VerbatimInput[numbers=none,fontsize=\scriptsize]{CodeSamples/formal/promela/qrcu.spin.col-ma.diff.lst}
\vspace*{-9pt}
\caption{\co{-DCOLLAPSE} 和 \co{-DMA=88} 的 Spin 输出差异}
\label{lst:formal:promela:Spin Output Diff of -DCOLLAPSE and -DMA=88}
\end{listing}

	作为间接证据，让我们比较
	使用 \co{-DCOLLAPSE} 和 \co{-DMA=88}
	（两个读者和三个更新者）运行的结果。
	这些运行的输出差异如
	\cref{lst:formal:promela:Spin Output Diff of -DCOLLAPSE and -DMA=88}所示。
	如你所见，它们在状态数（已存储和已匹配）上一致。
}\QuickQuizEnd

\begin{table*}
\rowcolors{6}{}{lightgray}
\renewcommand*{\arraystretch}{1.2}
\IfEbookSize{\scriptsize}{\footnotesize}
\centering
\OneColumnHSpace{-0.7in}%
\ebresizewidth{
\begin{tabular}{S[table-format = 1.0]S[table-format = 1.0]S[table-format = 9.0]
		S[table-format = 9.0]S[table-format = 2.0]S[table-format = 5.2]
		S[table-format = 4.2]S[table-format = 2.0]S[table-format = 4.2]
		S[table-format = 6.2]}
	\toprule
	\multicolumn{4}{r}{} & \multicolumn{3}{c}{\tco{-DCOLLAPSE}} &
					\multicolumn{3}{c}{\tco{-DMA=N}} \\
	\cmidrule(l){5-7} \cmidrule(l){8-10}
	\multicolumn{1}{r}{更新者} &
	    \multicolumn{1}{r}{读者} &
		\multicolumn{1}{r}{状态数} &
		    \multicolumn{1}{r}{达到的深度} &
			\multicolumn{1}{r}{\tco{-wN}} &
			    \multicolumn{1}{r}{内存 (MB)} &
				\multicolumn{1}{r}{运行时间 (s)} &
				    \multicolumn{1}{r}{\tco{N}} &
					\multicolumn{1}{r}{内存 (MB)} &
					    \multicolumn{1}{r}{运行时间 (s)} \\
	\cmidrule{1-4} \cmidrule(l){5-7} \cmidrule(l){8-10}
	1 & 1 &           376 &         95 & 12 &     0.10 & 0.00 &
		40 &    0.29 &      0.00 \\
	1 & 2 &         6 177 &        218 & 12 &     0.39 & 0.01 &
		47 &    0.59 &      0.02 \\
	1 & 3 &        99 728 &        385 & 16 &     4.60 & 0.14 &
		54 &    3.04 &      0.45 \\
        2 & 1 &        29 399 &        859 & 16 &     2.30 & 0.03 &
		55 &    0.70 &      0.13 \\
        2 & 2 &     1 071 181 &      2 352 & 20 &    49.24 & 1.45 &
		62 &    7.77 &      5.76 \\
        2 & 3 &    33 866 736 &     12 857 & 24 & 1 540.70 & 62.5 &
		69 &  111.66 &    326    \\
        3 & 1 &     2 749 453 &     53 809 & 21 &   125.25 & 4.01 &
		70 &   11.41 &     19.5  \\
        3 & 2 &   186 202 860 &    328 014 & 28 & 10 482.51 & 390 &
		77 &  222.26 &   2560    \\
	3 & 3 & 9 664 707 100 &  2 055 621 &    &          &      &
		84 & 5557.02 & 266000    \\
	\bottomrule
\end{tabular}
}
\caption{QRCU Spin 结果汇总}
\label{tab:formal:promela:QRCU Spin Result Summary}
\end{table*}

作为参考，\cref{tab:formal:promela:QRCU Spin Result Summary}
汇总了使用 \co{-DCOLLAPSE} 和 \co{-DMA=N}
编译器标志的 Spin 结果。
内存使用量是使用表中所示的最小足够搜索深度和
\co{-DMA=N} 参数获得的。
\co{-DCOLLAPSE} 运行的哈希表大小通过 \co{./pan} 的 \co{-wN}
选项进行调整，以避免在小状态空间上使用过多内存进行哈希。
因此内存使用量小于
\cref{tab:advsync:Memory Usage of QRCU Model}中所示的值，
后者的哈希表大小从默认的 \co{-w24} 开始。
运行时间来自一台 \Power{9} 服务器，这表明 \co{-DMA=N}
遭受的 CPU 开销比 \co{-DCOLLAPSE} 高约一个数量级，
但另一方面将内存开销减少了远超一个数量级。

目前为止一切顺利。
但再增加几个更新者或读者就会耗尽内存，
即使使用 \co{-DMA=N}。\footnote{
	或者，CPU 消耗会变得过大。}
那怎么办？
以下是一些可能的方法：

\begin{enumerate}
\item	看看较少数量的读者和更新者是否足以证明一般情况。
\item	手动构建正确性证明。
\item	使用更强大的工具。
\item	分而治之。
\end{enumerate}

以下各节讨论这些方法。

\subsubsection{到底需要多少读者和更新者？}
\label{sec:formal:How Many Readers and Updaters Are Really Needed?}

一种方法是仔细检查 \co{qrcu_updater()} 的 Promela 代码，
注意到唯一的全局状态变化发生在锁保护之下。
因此，在同一时刻只有一个更新者可能在修改
对读者或其他更新者可见的状态。
这意味着状态变化的任何序列都可以由单个更新者
串行执行，这是由于 Promela 进行完整的状态空间搜索。
因此，最多需要两个更新者：
一个用于修改状态，另一个用于将其搞乱。

读者的情况则不太明确，因为每个读者
仅仅执行一个单一的读端临界区然后退出。
可以论证有用的读者数量是受限的，
因为快速路径必定看到计数器中为零或者为一。
这是一个富有成效的研究方向，事实上，它导向了
下一节中描述的完整正确性证明。

\subsubsection{替代方法：
				     正确性证明}
\label{sec:formal:Alternative Approach: Proof of Correctness}

一个非正式证明~\cite{PaulMcKenney2007QRCUpatch}如下：

\begin{enumerate}
\item	如果 \co{synchronize_qrcu()} 过早退出，
	则根据定义，在 \co{synchronize_qrcu()} 的整个执行过程中
	必须至少有一个读者存在。
\item	对应该读者的计数器在此时间间隔内至少为 1。
\item	\co{synchronize_qrcu()} 代码强制
	至少一个计数器在任何时候至少为 1。
\item	以上两项意味着，如果对应该读者的计数器恰好为一，
	则另一个计数器必须大于或等于一。
	类似地，如果另一个计数器等于零，
	则对应读者的计数器必须大于或等于二。
\item	因此，在任何给定时间点，要么其中一个计数器至少为 2，
	要么两个计数器都至少为一。
\item	然而，\co{synchronize_qrcu()} 的快速路径代码
	在给定时间只能读取一个计数器。
	因此快速路径代码可能在获取第一个计数器时为零，
	但与计数器翻转竞争，使得第二个计数器被看到为一。
\item	在这样的竞态条件下，最多只能有一个读者持续存在，
	否则总和将为二或更大，
	这将导致更新者走慢速路径。
\item	但如果竞争发生在快速路径的第一次读取计数器时，
	然后在第二次读取时再次发生，
	则必须发生了两次计数器翻转。
\item	因为给定的更新者只翻转计数器一次，
	而且更新端锁阻止一对更新者并发翻转计数器，
	快速路径代码能够与翻转竞争两次的唯一方式是
	第一个更新者已经完成。
\item	但第一个更新者要等到所有先前存在的读者完成之后
	才会完成。
\item	因此，如果快速路径连续两次与计数器翻转竞争，
	所有先前存在的读者必定已经完成，
	所以走快速路径是安全的。
\end{enumerate}

当然，并非所有并行算法都有如此简单的证明。
在这种情况下，可能需要借助更强大的工具。

\subsubsection{替代方法：
				     更强大的工具}
\label{sec:formal:Alternative Approach: More Capable Tools}

虽然 Promela 和 Spin 十分有用，
但有更强大的工具可用，特别是用于验证硬件的工具。
这意味着如果能将你的算法翻译为硬件设计的 VHDL 语言——
对于非常底层的并行算法通常如此——
则可以将这些工具应用于你的代码
（例如，这曾被用于第一个实时 RCU 算法）。
然而，这些工具可能相当昂贵。

虽然商业化多处理器的出现可能最终会催生
具有强大状态空间缩减能力的自由软件模型检查器，
但这目前还没有什么大的帮助。

另外，Spin 有一些特性支持需要固定内存量的近似搜索，
然而，在验证并行算法时，我始终无法让自己信任这种近似方法。

另一种方法可能是分而治之。

\subsubsection{替代方法：
				     分而治之}
\label{sec:formal:Alternative Approach: Divide and Conquer}

通常可以将较大的并行算法分割成较小的片段，
然后分别证明它们。
例如，一个 100 亿状态的模型可能被分解为
一对 100,000 状态的模型。
采用这种方法不仅使 Promela 等工具更容易验证你的算法，
也会使你的算法更容易理解。

\subsubsection{QRCU 真的正确吗？}
\label{sec:formal:Is QRCU Really Correct?}

QRCU 真的正确吗？
我们有一个基于 Promela 的机械证明和一个手工证明，
二者都说它是正确的。
然而，\pplsur{Jade}{Alglave} 等人的一篇论文~\cite{JadeAlglave2013-cav}
却持相反意见（见论文第 12 页底部的 Section~5.1）。
到底怎么回事？

事实证明两者都是正确的！
当 QRCU 被添加到形式验证基准测试套件中时，
它的内存屏障被遗漏了，从而产生了一个有缺陷的 QRCU 版本。
所以这里真正的新闻是，许多形式验证工具
错误地证明了这个有缺陷的 QRCU 是正确的。
这就是为什么形式验证工具本身应该使用
被验证代码的注入缺陷版本来进行测试。
如果给定的工具无法找到注入的缺陷，那么它显然是不可信的。

\QuickQuiz{
	但不同的形式验证工具通常被设计为定位特定类别的缺陷。
	例如，很少有形式验证工具能发现规范中的错误。
	所以这个``显然不可信''的判断是不是有点苛刻？
}\QuickQuizAnswer{
	许多形式验证工具确实在某些方面是专门化的。
	例如，Promela 不处理现实的内存模型（memory model）
	（尽管可以将它们编程到
	\IX{Promela}~\cite{Desnoyers:2013:MSM:2506164.2506174} 中），
	\IXacr{cbmc}~\cite{EdmundClarke2004CBMC} 不检测概率性
	挂起和死锁（deadlock），
	\IX{Nidhugg}~\cite{CarlLeonardsson2014Nidhugg} 不检测
	涉及数据非确定性的缺陷。
	但这意味着这些工具不能被信任来发现
	它们不是为定位而设计的缺陷。

	因此，创建形式验证工具的人应该
	``在标签上说实话''，明确指出他们的工具
	能够和不能检测到哪些类别的缺陷。
	否则，当一个从业者第一次发现某个工具
	未能检测到一个缺陷时，该从业者很可能会
	对该工具进行极其严厉和极其公开的谴责。
	是的，是的，展示最好的一面确实有道理，
	但在没有适当免责声明的情况下走得太远，
	很容易触发负面反应的地雷，
	你的工具可能能也可能不能从中恢复。

	你已经被警告了！
}\QuickQuizEnd

因此，如果你确实打算使用 QRCU，请务必小心。
它的正确性证明本身可能正确也可能不正确。
这是形式验证不太可能完全取代测试的原因之一，
正如 Donald Knuth 很早以前就指出的那样。

\QuickQuiz{
	鉴于我们对此处描述的 QRCU 算法有两个独立的正确性证明，
	而且不正确性证明涵盖的是已知为不同的算法，
	为什么还有怀疑的余地？
}\QuickQuizAnswer{
	总是有怀疑的余地。
	在这种情况下，重要的是要记住这两个正确性证明
	先于现实世界内存模型的形式化，
	这就提出了这两个证明可能基于
	不正确的内存排序假设的可能性。
	此外，由于两个证明都是由同一个人构建的，
	它们很可能包含一个共同的错误。
	再说一次，总是有怀疑的余地。
}\QuickQuizEnd

% formal/dyntickrcu.tex
% mainfile: ../perfbook.tex
% SPDX-License-Identifier: CC-BY-SA-3.0

% Disable frame around VerbatimN in two-column layout
\IfTwoColumn{
\RecustomVerbatimEnvironment{VerbatimN}{Verbatim}%
{numbers=left,numbersep=5pt,xleftmargin=10pt,xrightmargin=0pt,frame=none}
\setlength{\lnlblraise}{0pt}
}{}

\subsection{Promela寓言：
			     {dynticks}与可抢占RCU}
\label{sec:formal:Promela Parable: dynticks and Preemptible RCU}

2008年初，一种可抢占的RCU变体被接受进入Linux主线内核，
以支持实时工作负载。
该变体类似于自2005年8月以来
\rt\ 补丁集~\cite{IngoMolnar05a}中的RCU实现。
实时工作负载需要可抢占RCU，因为早期的RCU实现在整个RCU
读侧临界区内禁用了抢占，导致过高的实时延迟。

然而，早期\rt\ 实现的一个缺点是每个\IX{grace period}
都需要在每个CPU上执行工作，即使该CPU处于低功耗的
``dynticks-idle''状态，因而无法执行RCU读侧临界区。
dynticks-idle状态背后的理念是，空闲的CPU应当被完全断电以
节省能源。
简而言之，可抢占RCU可能会使近期Linux内核中一项
有价值的节能特性失效。
尽管Josh Triplett和Paul McKenney曾讨论过一些方法，
允许CPU在整个RCU宽限期内保持低功耗状态
（从而保留Linux内核的节能能力），但直到Steve Rostedt
在\rt\ 补丁集中将一个新的dyntick实现与可抢占RCU集成后，
事情才到了非解决不可的地步。

这种组合导致Steve的一台系统在启动时挂起，因此在2007年十月，
Paul编写了一个对可抢占RCU宽限期处理的dynticks友好修改。
Steve编写了\co{rcu_irq_enter()}和\co{rcu_irq_exit()}
接口，从\co{irq_enter()}和\co{irq_exit()}中断
入口/出口函数中调用。
这些\co{rcu_irq_enter()}和\co{rcu_irq_exit()}
函数是必需的，以允许RCU可靠地处理dynticks-idle CPU
被瞬间唤醒以执行包含RCU读侧临界区的中断处理程序的情况。
有了这些修改，Steve的系统能够可靠启动。
但Paul假定我们不可能在第一次就把代码写对了，
因此继续定期检查代码。

Paul从2007年10月到2008年2月反复审查代码，
几乎每次都能发现至少一个bug。
有一次，Paul甚至在意识到一个bug是虚幻的之前，
就已经编写并测试了修复程序。
实际上在所有情况下，这些``bug''都被证明是虚幻的。

二月底时，Paul对这个游戏已经有点疲倦了。
因此他决定借助Promela和Spin的帮助。
以下内容呈现了一系列七个逐步增加复杂度的Promela模型，
最后一个通过了验证，状态空间消耗了约40\,GB的主内存。

更重要的是，Promela和Spin确实为我找到了一个非常微妙的bug！

\QuickQuiz{
	好吧，那真是太好了！
	现在，如果我碰巧没有一台配备40\,GB主内存的机器，
	我该怎么办？？？
}\QuickQuizAnswer{
	别急，这个问题有不少合法的答案：
	\begin{enumerate}
	\item	尝试编译器标志\co{-DCOLLAPSE}和\co{-DMA=N}
		以减少内存消耗。
		参见\cref{sec:formal:Running the QRCU Example}。
	\item	进一步优化模型，减少其内存消耗。
	\item	进行纸笔证明，也许可以从Linux内核代码中的
		注释开始。
	\item	设计细致的压力测试，虽然它们不能证明代码正确，
		但可以发现隐藏的bug。
	\item	目前有一些趋势是开发能在较小机器集群上进行
		模型检查的工具。
		但请注意，由于Paul偶尔能使用一些大型机器，
		我们实际上并没有使用过这类工具。
	\item	等待经济适用系统的内存容量扩展到能适应你
		的问题。
	\item	使用云计算服务之一来短期租用一台大型系统。
	\end{enumerate}
}\QuickQuizEnd

更好的做法是设计出一个更简单、更快速且状态空间更小的算法。
最好的情况是算法如此简单，其正确性对于一般观察者来说
显而易见！

\Crefrange{sec:formal:Introduction to Preemptible RCU and dynticks}
{sec:formal:Grace-Period Interface}
概述了可抢占RCU的dynticks接口，
随后
\cref{sec:formal:Validating Preemptible RCU and dynticks}
讨论了该接口的验证。

\subsubsection{可抢占RCU与dynticks简介}
\label{sec:formal:Introduction to Preemptible RCU and dynticks}

每CPU变量\co{dynticks_progress_counter}是dynticks与
可抢占RCU之间接口的核心。
当对应的CPU处于dynticks-idle模式时，该变量的值为偶数，
否则为奇数。
CPU退出dynticks-idle模式有以下三个原因：

\begin{enumerate}
\item	开始运行一个任务，
\item	进入可能嵌套的一组中断处理程序的最外层时，以及
\item	进入\IXacr{nmi}处理程序时。
\end{enumerate}

可抢占RCU的宽限期机制对\co{dynticks_progress_counter}
变量的值进行采样，以确定何时可以安全地忽略处于
dynticks-idle状态的CPU。

以下三节概述了任务接口、中断/NMI接口，
以及Linux内核v2.6.25-rc4中宽限期机制对
\co{dynticks_progress_counter}变量的使用。

\subsubsection{任务接口}
\label{sec:formal:Task Interface}

当某个CPU因为没有更多任务可运行而进入dynticks-idle模式时，
它会调用\co{rcu_enter_nohz()}：

\begin{VerbatimN}
static inline void rcu_enter_nohz(void)
{
	mb();
	__get_cpu_var(dynticks_progress_counter)++;
	WARN_ON(__get_cpu_var(dynticks_progress_counter) &
	        0x1);
}
\end{VerbatimN}

这个函数简单地递增\co{dynticks_progress_counter}并检查结果
是否为偶数，但在此之前先执行一个内存屏障，以确保任何其他
CPU在看到\co{dynticks_progress_counter}的新值时，也能看到
之前所有RCU读侧临界区的完成。

类似地，当处于dynticks-idle模式的CPU准备开始执行一个
新的可运行任务时，它会调用\co{rcu_exit_nohz()}：

\begin{VerbatimN}
static inline void rcu_exit_nohz(void)
{
	__get_cpu_var(dynticks_progress_counter)++;
	mb();
	WARN_ON(!(__get_cpu_var(dynticks_progress_counter) &
	          0x1));
}
\end{VerbatimN}

这个函数同样递增\co{dynticks_progress_counter}，
但在其后跟随一个内存屏障，以确保如果任何其他CPU
看到后续RCU读侧临界区的结果，那么该CPU也将看到
\co{dynticks_progress_counter}递增后的值。
最后，\co{rcu_exit_nohz()}检查递增的结果是否为奇数值。

\co{rcu_enter_nohz()}和\co{rcu_exit_nohz()}
函数处理CPU因任务执行而进入和退出dynticks-idle模式的情况，
但不处理中断，中断将在下一节中介绍。

\subsubsection{中断接口}
\label{sec:formal:Interrupt Interface}

\co{rcu_irq_enter()}和\co{rcu_irq_exit()}
函数分别处理中断/NMI的进入和退出。
当然，嵌套中断也必须被正确处理。
嵌套中断的可能性通过第二个每CPU变量\co{rcu_update_flag}
来处理，该变量在进入中断或NMI处理程序时（在
\co{rcu_irq_enter()}中）递增，
在退出时（在\co{rcu_irq_exit()}中）递减。
此外，已有的\co{in_interrupt()}原语用于区分最外层
中断和嵌套中断/NMI\@。

中断进入由如下所示的\co{rcu_irq_enter()}处理：

\begin{fcvlabel}[ln:formal:dyntickrcu:rcu_irq_enter]
\begin{VerbatimN}[commandchars=\\\[\]]
void rcu_irq_enter(void)
{
	int cpu = smp_processor_id();	\lnlbl[fetch]

	if (per_cpu(rcu_update_flag, cpu))	\lnlbl[inc:b]
		per_cpu(rcu_update_flag, cpu)++; \lnlbl[inc:e]
	if (!in_interrupt() &&			\lnlbl[chk_lv:b]
	    (per_cpu(dynticks_progress_counter,
	             cpu) & 0x1) == 0) {	\lnlbl[chk_lv:e]
		per_cpu(dynticks_progress_counter, cpu)++; \lnlbl[inc_cnt]
		smp_mb();			\lnlbl[mb]
		per_cpu(rcu_update_flag, cpu)++;\lnlbl[inc_flg]
	}
}
\end{VerbatimN}
\end{fcvlabel}

\begin{fcvref}[ln:formal:dyntickrcu:rcu_irq_enter]
\Clnref{fetch}获取当前CPU的编号，而\clnref{inc:b,inc:e}
在\co{rcu_update_flag}嵌套计数器已经非零时将其递增。
\Clnrefrange{chk_lv:b}{chk_lv:e}检查我们是否处于中断的
最外层，如果是，还检查\co{dynticks_progress_counter}
是否需要递增。
如果需要，\clnref{inc_cnt}递增\co{dynticks_progress_counter}，
\clnref{mb}执行内存屏障，\clnref{inc_flg}递增
\co{rcu_update_flag}。
与\co{rcu_exit_nohz()}一样，内存屏障确保任何其他CPU
在看到中断处理程序中RCU读侧临界区的效果时（在
\co{rcu_irq_enter()}调用之后），也能看到
\co{dynticks_progress_counter}的递增。
\end{fcvref}

\QuickQuizSeries{%
\QuickQuizB{
	为什么不简单地递增\co{rcu_update_flag}，然后仅在
	\co{rcu_update_flag}的旧值为零时才递增
	\co{dynticks_progress_counter}？？？
}\QuickQuizAnswerB{
	这在存在NMI的情况下会失败。
	要理解这一点，假设在\co{rcu_irq_enter()}递增了
	\co{rcu_update_flag}之后、但在递增
	\co{dynticks_progress_counter}之前收到了一个NMI。
	由NMI调用的\co{rcu_irq_enter()}实例将看到
	\co{rcu_update_flag}的原始值为非零，因此会放弃递增
	\co{dynticks_progress_counter}。
	这将使RCU宽限期机制无法知道NMI处理程序正在该CPU上
	执行，导致NMI处理程序中的任何RCU读侧临界区失去RCU保护。

	NMI处理程序的存在——按定义不可被屏蔽——确实使这段
	代码变得复杂。
}\QuickQuizEndB
%
\QuickQuizE{
	\begin{fcvref}[ln:formal:dyntickrcu:rcu_irq_enter]
	但是如果\clnref{chk_lv:b}发现我们处于最外层中断，
	我们不是\emph{总是}需要递增
	\co{dynticks_progress_counter}吗？
	\end{fcvref}
}\QuickQuizAnswerE{
	如果我们中断的是一个正在运行的任务则不需要！
	在那种情况下，\co{dynticks_progress_counter}已经被
	\co{rcu_exit_nohz()}递增过了，不需要再次递增。
}\QuickQuizEndE
}

中断退出由\co{rcu_irq_exit()}以类似方式处理：

\begin{fcvlabel}[ln:formal:dyntickrcu:rcu_irq_exit]
\begin{VerbatimN}[commandchars=\\\[\]]
void rcu_irq_exit(void)
{
	int cpu = smp_processor_id();	\lnlbl[fetch]

	if (per_cpu(rcu_update_flag, cpu)) { \lnlbl[chk_flg]
		if (--per_cpu(rcu_update_flag, cpu)) \lnlbl[dec_flg]
			return;
		WARN_ON(in_interrupt());	\lnlbl[verify]
		smp_mb();			\lnlbl[mb]
		per_cpu(dynticks_progress_counter, cpu)++; \lnlbl[inc_cnt]
		WARN_ON(per_cpu(dynticks_progress_counter, \lnlbl[vrf_even:b]
		                cpu) & 0x1); \lnlbl[vrf_even:e]
	}
}
\end{VerbatimN}
\end{fcvlabel}

\begin{fcvref}[ln:formal:dyntickrcu:rcu_irq_exit]
\Clnref{fetch}与之前一样获取当前CPU的编号。
\Clnref{chk_flg}检查\co{rcu_update_flag}是否非零，
如果为零则立即返回（通过从函数末尾落出）。
否则，\clnrefthro{dec_flg}{vrf_even:e}开始发挥作用。
\Clnref{dec_flg}递减\co{rcu_update_flag}，如果结果不为零
则返回。
\Clnref{verify}验证我们确实正在离开嵌套中断的最外层，
\clnref{mb}执行内存屏障，
\clnref{inc_cnt}递增\co{dynticks_progress_counter}，
\clnref{vrf_even:b,vrf_even:e}验证该变量现在是偶数。
与\co{rcu_enter_nohz()}一样，内存屏障确保任何其他CPU
在看到\co{dynticks_progress_counter}的递增时，
也能看到中断处理程序中RCU读侧临界区的效果
（在\co{rcu_irq_exit()}调用之前）。
\end{fcvref}

以上两节描述了在任务、中断和NMI进入和退出dynticks-idle
模式期间，\co{dynticks_progress_counter}变量是如何维护的。
下一节描述可抢占RCU的宽限期机制如何使用该变量。

\subsubsection{宽限期接口}
\label{sec:formal:Grace-Period Interface}

在\cref{fig:formal:Preemptible RCU State Machine}所示的
四个可抢占RCU宽限期状态中，
只有\co{rcu_try_flip_waitack_state}
和\co{rcu_try_flip_waitmb_state}状态需要等待
其他CPU响应。

\begin{figure}
\centering
\resizebox{2.5in}{!}{\includegraphics{formal/RCUpreemptStates}}
\caption{可抢占RCU状态机}
\label{fig:formal:Preemptible RCU State Machine}
\end{figure}

当然，如果给定的CPU处于dynticks-idle状态，我们不应该
等待它。
因此，在进入这两个状态之一之前，前一个状态会对每个CPU的
\co{dynticks_progress_counter}变量进行快照，并将快照
存放在另一个每CPU变量\co{rcu_dyntick_snapshot}中。
这通过调用如下所示的
\co{dyntick_save_progress_counter()}来完成：

\begin{VerbatimN}
static void dyntick_save_progress_counter(int cpu)
{
	per_cpu(rcu_dyntick_snapshot, cpu) =
		per_cpu(dynticks_progress_counter, cpu);
}
\end{VerbatimN}

\co{rcu_try_flip_waitack_state}状态调用如下所示的
\co{rcu_try_flip_waitack_needed()}：

\begin{fcvlabel}[ln:formal:dyntickrcu:rcu_try_flip_waitack_needed]
\begin{VerbatimN}[samepage=true,commandchars=\\\[\]]
static inline int
rcu_try_flip_waitack_needed(int cpu)
{
	long curr;
	long snap;

	curr = per_cpu(dynticks_progress_counter, cpu); \lnlbl[curr]
	snap = per_cpu(rcu_dyntick_snapshot, cpu); \lnlbl[snap]
	smp_mb();				\lnlbl[mb]
	if ((curr == snap) && ((curr & 0x1) == 0)) \lnlbl[chk_remain]
		return 0;			\lnlbl[ret_0_r]
	if ((curr - snap) > 2 || (snap & 0x1) == 0) \lnlbl[chk_idle]
		return 0;			\lnlbl[ret_0_i]
	return 1;				\lnlbl[ret_1]
}
\end{VerbatimN}
\end{fcvlabel}

\begin{fcvref}[ln:formal:dyntickrcu:rcu_try_flip_waitack_needed]
\Clnref{curr,snap}分别获取\co{dynticks_progress_counter}的
当前版本和快照版本。
\clnref{mb}上的内存屏障确保后续
\co{rcu_try_flip_waitzero_state}中的计数器检查
在这些计数器获取之后执行。
\Clnref{chk_remain,ret_0_r}在该CPU自快照获取以来一直
保持在dynticks-idle状态时返回零（意味着不需要与
指定CPU通信）。
类似地，\clnref{chk_idle,ret_0_i}在该CPU最初处于
dynticks-idle状态或已经完整地经历了一次
dynticks-idle状态时返回零。
在这两种情况下，CPU不可能保留宽限期计数器的旧值。
如果这两个条件都不满足，\clnref{ret_1}返回一，
意味着该CPU需要显式响应。
\end{fcvref}

对于\co{rcu_try_flip_waitmb_state}状态，它调用如下所示的
\co{rcu_try_flip_waitmb_needed()}：

\begin{fcvlabel}[ln:formal:dyntickrcu:rcu_try_flip_waitmb_needed]
\begin{VerbatimN}[commandchars=\\\[\]]
static inline int
rcu_try_flip_waitmb_needed(int cpu)
{
	long curr;
	long snap;

	curr = per_cpu(dynticks_progress_counter, cpu);
	snap = per_cpu(rcu_dyntick_snapshot, cpu);
	smp_mb();
	if ((curr == snap) && ((curr & 0x1) == 0))
		return 0;
	if (curr != snap)		\lnlbl[chk_to_from]
		return 0;		\lnlbl[ret_0]
	return 1;
}
\end{VerbatimN}
\end{fcvlabel}

\begin{fcvref}[ln:formal:dyntickrcu:rcu_try_flip_waitmb_needed]
这与\co{rcu_try_flip_waitack_needed()}非常相似，
不同之处在于\clnref{chk_to_from,ret_0}，因为任何进入或
离开dynticks-idle状态的转换都会执行
\co{rcu_try_flip_waitmb_state}状态所需的内存屏障。
\end{fcvref}

我们现在已经看到了RCU与dynticks-idle状态之间接口涉及的
所有代码。
下一节将构建用于验证这些代码的Promela模型。

\QuickQuiz{
	你能在本节的代码中发现任何bug吗？
}\QuickQuizAnswer{
	阅读下一节来看看你是否正确。
}\QuickQuizEnd

\subsection{验证可抢占RCU和dynticks}
\label{sec:formal:Validating Preemptible RCU and dynticks}

本节逐步开发dynticks与RCU之间接口的Promela模型，
\crefrange{sec:formal:Basic Model}{sec:formal:Validating NMI Handlers}
中的每一节都展示了一个步骤，从进程级代码开始，
逐步添加断言、中断，最后是NMI。

\Cref{sec:formal:Lessons (Re)Learned}列出了
在这一过程中（重新）学到的经验教训，
\crefrange{sec:formal:Simplicity Avoids Formal Verification}{sec:formal:Discussion}
则展示了一个更简单的RCU dynticks问题解决方案。

\subsubsection{基本模型}
\label{sec:formal:Basic Model}

本节将进程级的dynticks进入/退出代码和宽限期处理翻译为
Promela~\cite{Holzmann03a}。
我们从2.6.25-rc4内核的\co{rcu_exit_nohz()}和
\co{rcu_enter_nohz()}开始，
将它们放入一个Promela进程中，以循环方式建模退出和
进入dynticks-idle模式，如下所示：

\input{CodeSamples/formal/promela/dyntick/dyntickRCU-base@dyntick_nohz.fcv}

\begin{fcvref}[ln:formal:promela:dyntick:dyntickRCU-base:dyntick_nohz]
\Clnref{do,od}定义了一个循环。
\Clnref{break}在循环计数器~\co{i}超过限制
\co{MAX_DYNTICK_LOOP_NOHZ}时退出循环。
\Clnref{kick_loop}告诉循环构造在每次循环迭代中执行
\clnrefrange{ex_inc:b}{inc_i}。
因为\clnref{break,kick_loop}上的条件是互斥的，
Promela通常的随机选择真条件的行为被禁用。
\Clnref{ex_inc:b,ex_inc:e}建模了\co{rcu_exit_nohz()}对
\co{dynticks_progress_counter}的非原子递增，而
\clnref{ex_assert}建模了\co{WARN_ON()}。
\co{atomic}构造只是为了减少Promela的状态空间，
因为\co{WARN_ON()}严格来说不是算法的一部分。
\Clnrefrange{ent_inc:b}{ent_inc:e}类似地建模了
\co{rcu_enter_nohz()}的递增和\co{WARN_ON()}。
最后，\clnref{inc_i}递增循环计数器。
\end{fcvref}

因此，循环的每次迭代建模了一个CPU退出dynticks-idle模式
（例如，开始执行一个任务），然后重新进入dynticks-idle模式
（例如，该任务阻塞）。

\QuickQuizSeries{%
\QuickQuizB{
	为什么在Promela中不对\co{rcu_exit_nohz()}和
	\co{rcu_enter_nohz()}中的内存屏障进行建模？
}\QuickQuizAnswerB{
	Promela假设顺序一致性，因此没有必要对内存屏障建模。
	实际上，人们反而必须显式地对缺少内存屏障的情况建模，
	例如，如
	\cref{lst:formal:QRCU Unordered Summation}
	（\cpageref{lst:formal:QRCU Unordered Summation}）
	所示。
}\QuickQuizEndB
%
\QuickQuizE{
	先建模\co{rcu_exit_nohz()}再建模
	\co{rcu_enter_nohz()}是不是有点奇怪？
	先建模进入再建模退出不是更自然吗？
}\QuickQuizAnswerE{
	那样可能确实更自然，但我们在后面要添加的活性检查中
	需要这种特定的顺序。
}\QuickQuizEndE
}

下一步是建模RCU宽限期处理的接口。
为此，我们需要建模
\co{dyntick_save_progress_counter()}、
\co{rcu_try_flip_waitack_needed()}、
\co{rcu_try_flip_waitmb_needed()}，
以及\co{rcu_try_flip_waitack()}和
\co{rcu_try_flip_waitmb()}的部分内容，
均来自2.6.25-rc4内核。
下面的\co{grace_period()} Promela进程建模了这些函数
在可抢占RCU宽限期处理的单次遍历中的调用方式。

\input{CodeSamples/formal/promela/dyntick/dyntickRCU-base@grace_period.fcv}

\begin{fcvref}[ln:formal:promela:dyntick:dyntickRCU-base:grace_period]
\Clnrefrange{print:b}{print:e}打印循环限制
（但只在出错时打印到``.trail''文件中），
并建模了\co{rcu_try_flip_idle()}中的一行代码及其对
\co{dyntick_save_progress_counter()}的调用，
该调用对当前CPU的\co{dynticks_progress_counter}变量
进行快照。
这两行原子地执行以减少状态空间。

\Clnrefrange{do1}{od1}建模了\co{rcu_try_flip_waitack()}
中的相关代码及其对
\co{rcu_try_flip_waitack_needed()}的调用。
该循环建模了宽限期状态机等待每个CPU的计数器翻转确认，
但仅涉及与dynticks-idle CPU交互的部分。

\Clnref{snap}建模了\co{rcu_try_flip_waitzero()}中的
一行及其对\co{dyntick_save_progress_counter()}的调用，
再次对CPU的\co{dynticks_progress_counter}变量进行快照。

最后，\clnrefrange{do2}{od2}建模了
\co{rcu_try_flip_waitack()}中的相关代码及其对
\co{rcu_try_flip_waitack_needed()}的调用。
该循环建模了宽限期状态机等待每个CPU执行内存屏障，
但同样仅涉及与dynticks-idle CPU交互的部分。
\end{fcvref}

\QuickQuiz{
	等一下！
	在Linux内核中，\co{dynticks_progress_counter}和
	\co{rcu_dyntick_snapshot}都是每CPU变量。
	那为什么它们在这里被建模为单个全局变量？
}\QuickQuizAnswer{
	因为宽限期代码分别处理每个CPU的
	\co{dynticks_progress_counter}和
	\co{rcu_dyntick_snapshot}变量，
	我们可以将状态压缩到单个CPU上。
	如果宽限期代码在特定CPU上针对特定值做特殊处理，
	那么我们确实需要建模多个CPU。
	但幸运的是，我们可以安全地限制在两个CPU上：
	一个运行宽限期处理，另一个进入和离开dynticks-idle模式。
}\QuickQuizEnd

得到的模型（\path{dyntickRCU-base.spin}），
使用\path{runspin.sh}脚本运行时，
生成了691个状态并且无错误通过，
鉴于该模型完全缺少能够发现故障的断言，
这一点丝毫不令人意外。
因此，下一节添加安全性断言。

\subsubsection{验证安全性}
\label{sec:formal:Validating Safety}

一个安全的RCU实现绝不能允许宽限期在任何在宽限期开始
之前启动的RCU读者完成之前结束。
这通过一个\co{grace_period_state}变量来建模，
该变量可以取以下三种状态：

\input{CodeSamples/formal/promela/dyntick/dyntickRCU-base-s@grace_period_state.fcv}

\co{grace_period()}进程在经过宽限期各阶段时设置该变量，
如下所示：

\input{CodeSamples/formal/promela/dyntick/dyntickRCU-base-s@grace_period.fcv}

\begin{fcvref}[ln:formal:promela:dyntick:dyntickRCU-base-s:grace_period]
\Clnref{upd_gps1,upd_gps2,upd_gps3,upd_gps4,upd_gps5,upd_gps6}
更新该变量（在可行的地方与算法操作原子地组合），
以允许\co{dyntick_nohz()}进程验证基本的RCU安全性属性。
此验证的形式是断言\co{grace_period_state}变量的值
不能在RCU读者可能持续存在的时间段内从
\co{GP_IDLE}跳到\co{GP_DONE}。
\end{fcvref}

\QuickQuiz{
	\begin{fcvref}[ln:formal:promela:dyntick:dyntickRCU-base-s:grace_period]
	既然在\clnref{upd_gps3,upd_gps4}上有一对背靠背的
	\co{grace_period_state}修改，
	我们如何确保\clnref{upd_gps3}的修改不会丢失？
	\end{fcvref}
}\QuickQuizAnswer{
	回想一下Promela和Spin会追踪每一种可能的状态变化序列。
	因此，时序是无关紧要的：
	Promela/Spin会很乐意在这两条语句之间插入模型的
	所有剩余部分，除非某个状态变量明确禁止这样做。
}\QuickQuizEnd

\co{dyntick_nohz()} Promela进程实现此验证如下所示：

\input{CodeSamples/formal/promela/dyntick/dyntickRCU-base-s@dyntick_nohz.fcv}

\begin{fcvref}[ln:formal:promela:dyntick:dyntickRCU-base-s:dyntick_nohz]
\Clnref{new_flg}在任务执行开始时如果
\co{grace_period_state}变量的值为\co{GP_IDLE}，
则设置一个新的\co{old_gp_idle}标志，
而\clnref{assert:b,assert:e}处的断言在
\co{grace_period_state}变量在任务执行期间
推进到\co{GP_DONE}时触发，
这是非法的，因为单个RCU读侧临界区可能跨越整个
中间时间段。
\end{fcvref}

得到的模型（\path{dyntickRCU-base-s.spin}），
使用\path{runspin.sh}脚本运行时，
生成了964个状态并无错误通过，这令人放心。
尽管如此，虽然安全性至关重要，避免宽限期无限停滞也同样
重要。因此下一节将介绍活性验证。

\subsubsection{验证活性}
\label{sec:formal:Validating Liveness}

\begin{fcvref}[ln:formal:promela:dyntick:dyntickRCU-base-sl-busted:dyntick_nohz]
虽然活性可能很难证明，但这里有一个简单的技巧适用。
第一步是让\co{dyntick_nohz()}通过一个
\co{dyntick_nohz_done}变量来指示它已完成，
如下面的\clnref{done}所示：
\end{fcvref}

\input{CodeSamples/formal/promela/dyntick/dyntickRCU-base-sl-busted@dyntick_nohz.fcv}

有了这个变量，我们可以在\co{grace_period()}中添加断言
来检查是否存在不必要的阻塞，如下所示：

\input{CodeSamples/formal/promela/dyntick/dyntickRCU-base-sl-busted@grace_period.fcv}

\begin{fcvref}[ln:formal:promela:dyntick:dyntickRCU-base-sl-busted:grace_period]
我们在\clnref{shex}添加了\co{shouldexit}变量，
在\clnref{init_shex}将其初始化为零。
\Clnref{assert_shex}断言\co{shouldexit}未被设置，
而\clnref{set_shex}将\co{shouldexit}设置为
\co{dyntick_nohz()}维护的\co{dyntick_nohz_done}变量。
因此，如果我们在\co{dyntick_nohz()}执行完成后
尝试在等待计数器翻转确认的循环中再次迭代，
该断言就会触发。
毕竟，如果\co{dyntick_nohz()}已完成，就不会再有状态
变化来强制我们退出循环，因此在这种状态下迭代两次
意味着无限循环，进而意味着宽限期永远不会结束。

\Clnref{init_shex2,assert_shex2,set_shex2}对第二个
（内存屏障）循环以类似方式操作。

然而，运行该模型（\path{dyntickRCU-base-sl-busted.spin}）
会导致失败，因为\clnref{chk_2}检查了错误的变量
是否为偶数。
失败时，\co{spin}会写出一个``trail''文件
（\path{dyntickRCU-base-sl-busted.spin.trail}），
记录导致失败的状态序列。
使用``\co{spin -t -p -g -l\ }%
\path{dyntickRCU-base-sl-busted.spin}''
命令可以让\co{spin}重新追踪这个状态序列，
打印执行的语句和变量的值
（\path{dyntickRCU-base-sl-busted.spin.trail.txt}）。
请注意，由于\co{spin}将两个函数放在单个文件中，
行号与上面的清单不匹配。
但是，行号\emph{确实}与完整模型
（\path{dyntickRCU-base-sl-busted.spin}）匹配。

我们看到\co{dyntick_nohz()}进程在步骤34完成
（搜索``34:''），但\co{grace_period()}进程仍然
未能退出循环。
\co{curr}的值为\co{6}（见步骤35），
\co{snap}的值为\co{5}（见步骤17）。
因此\clnref{chk_1}上的第一个条件不成立，因为
\qco{curr != snap}，而\clnref{chk_2}上的第二个条件
也不成立，因为\co{snap}是奇数且\co{curr}仅比
\co{snap}大一。
\end{fcvref}

所以这两个条件中必有一个是不正确的。
参考\co{rcu_try_flip_waitack_needed()}中
第一个条件的注释块：

\begin{quote}
	如果CPU在整个期间保持在dynticks模式且没有接收任何
	中断、NMI、SMI或其他，那么它不可能处于
	\co{rcu_read_lock()}的中间，所以它执行的下一个
	\co{rcu_read_lock()}必须使用计数器的新值。
	因此我们可以安全地假装该CPU已经确认了计数器。
\end{quote}

第一个条件确实与此匹配，因为如果\qco{curr == snap}
且\co{curr}为偶数，那么对应的CPU在整个期间一直处于
dynticks-idle模式，符合要求。
那么我们来看看第二个条件的注释块：

\begin{quote}
	如果CPU在没有活跃irq处理程序的情况下经历了或进入了
	dynticks空闲阶段，那么如上所述，我们可以安全地假装
	该CPU已经确认了计数器。
\end{quote}

条件的第一部分是正确的，因为如果\co{curr}和\co{snap}
相差二，则中间至少有一个偶数，对应于完整地经历了一个
dynticks-idle阶段。
然而，条件的第二部分对应的是\emph{开始}于dynticks-idle
模式，而不是\emph{结束}于该模式。
因此我们需要测试的是\co{curr}而非\co{snap}是否为偶数。

修正后的C代码如下：

\begin{fcvlabel}[ln:formal:dyntickrcu:rcu_try_flip_waitack_needed_fixed]
\begin{VerbatimN}[commandchars=\\\[\]]
static inline int
rcu_try_flip_waitack_needed(int cpu)
{
	long curr;
	long snap;

	curr = per_cpu(dynticks_progress_counter, cpu);
	snap = per_cpu(rcu_dyntick_snapshot, cpu);
	smp_mb();
	if ((curr == snap) && ((curr & 0x1) == 0)) \lnlbl[if:b]
		return 0;
	if ((curr - snap) > 2 || (curr & 0x1) == 0)
		return 0;			\lnlbl[if:e]
	return 1;
}
\end{VerbatimN}
\end{fcvlabel}

\begin{fcvref}[ln:formal:dyntickrcu:rcu_try_flip_waitack_needed_fixed]
\Clnrefrange{if:b}{if:e}现在可以合并和简化，
得到如下结果。
类似的简化也可以应用于\co{rcu_try_flip_waitmb_needed()}。
\end{fcvref}

\begin{VerbatimN}[commandchars=\\\[\]]
static inline int
rcu_try_flip_waitack_needed(int cpu)
{
	long curr;
	long snap;

	curr = per_cpu(dynticks_progress_counter, cpu);
	snap = per_cpu(rcu_dyntick_snapshot, cpu);
	smp_mb();
	if ((curr - snap) >= 2 || (curr & 0x1) == 0)
		return 0;
	return 1;
}
\end{VerbatimN}

在模型（\path{dyntickRCU-base-sl.spin}）中进行相应修正后，
得到一个包含661个状态的正确验证，无错误通过。
然而，值得注意的是，活性验证的第一个版本未能捕获这个bug，
原因是活性验证本身存在一个bug。
这个活性验证的bug是通过在\co{grace_period()}进程中
插入一个无限循环，并注意到活性验证代码未能检测到
该问题而发现的！

我们现在已经成功验证了安全性和活性条件，但仅限于进程运行
和阻塞的情况。
我们还需要处理中断，这是下一节要讨论的任务。

\subsubsection{中断}
\label{sec:formal:Interrupts}

在Promela中有两种方式来建模中断：
\begin{enumerate}
\item	使用C预处理器技巧在\co{dynticks_nohz()}进程的
	每对相邻语句之间插入中断处理程序，或者
\item	用一个单独的进程来建模中断处理程序。
\end{enumerate}

经过一番思考，第二种方法的状态空间会更小，
尽管它要求中断处理程序相对于\co{dynticks_nohz()}进程
原子执行，但不需要相对于\co{grace_period()}进程
原子执行。

幸运的是，Promela允许从原子语句中分支跳出。
这个技巧允许我们让中断处理程序设置一个标志，
并重新编写\co{dynticks_nohz()}以原子地检查该标志，
仅在标志未设置时才执行。
这可以通过一个接受标签和Promela语句的C预处理器宏来实现，
如下所示：

\input{CodeSamples/formal/promela/dyntick/dyntickRCU-irqnn-ssl@EXECUTE_MAINLINE.fcv}

可以如下使用该宏：

\begin{VerbatimU}
EXECUTE_MAINLINE(stmt1,
                 tmp = dynticks_progress_counter)
\end{VerbatimU}

\begin{fcvref}[ln:formal:promela:dyntick:dyntickRCU-irqnn-ssl:EXECUTE_MAINLINE]
宏的\clnref{label}创建指定的语句标签。
\Clnrefrange{atm:b}{atm:e}是一个原子块，测试
\co{in_dyntick_irq}变量，如果该变量被设置
（表示中断处理程序处于活动状态），则从原子块跳出
回到标签处。
否则，\clnref{else}执行指定的语句。
总体效果是，主线执行在中断活动时暂停，这正是所需要的。
\end{fcvref}

\subsubsection{验证中断处理程序}
\label{sec:formal:Validating Interrupt Handlers}

第一步是将\co{dyntick_nohz()}转换为
\co{EXECUTE_MAINLINE()}形式，如下所示：

\input{CodeSamples/formal/promela/dyntick/dyntickRCU-irqnn-ssl@dyntick_nohz.fcv}

\begin{fcvref}[ln:formal:promela:dyntick:dyntickRCU-irqnn-ssl:dyntick_nohz]
需要注意的是，当一组语句被传递给\co{EXECUTE_MAINLINE()}时，
如\clnrefrange{stmt2:b}{stmt2:e}，该组中的所有语句都原子地执行。
\end{fcvref}

\QuickQuizSeries{%
\QuickQuizB{
	但是如果你需要单个\co{EXECUTE_MAINLINE()}组中的
	语句非原子地执行，你会怎么做？
}\QuickQuizAnswerB{
	最简单的做法是将每个这样的语句放入自己的
	\co{EXECUTE_MAINLINE()}语句中。
}\QuickQuizEndB
%
\QuickQuizE{
	但是如果\co{dynticks_nohz()}进程有带条件的
	\qco{if}或\qco{do}语句，而这些构造的语句体
	需要非原子地执行呢？
}\QuickQuizAnswerE{
	一种方法，正如我们将在后面的章节中看到的，
	是使用显式标签和\qco{goto}语句。
	例如，以下构造：

\begin{VerbatimU}
	if
	:: i == 0 -> a = -1;
	:: else -> a = -2;
	fi;
\end{VerbatimU}
%
	可以建模为类似如下的形式：

\begin{VerbatimU}
	EXECUTE_MAINLINE(stmt1,
	                 if
	                 :: i == 0 -> goto stmt1_then;
	                 :: else -> goto stmt1_else;
	                 fi)
	stmt1_then: skip;
	EXECUTE_MAINLINE(stmt1_then1, a = -1; goto stmt1_end)
	stmt1_else: skip;
	EXECUTE_MAINLINE(stmt1_then1, a = -2)
	stmt1_end: skip;
\end{VerbatimU}

	然而，在\qco{if}语句的情况下，该宏的帮助并不明显，
	因此这类情况将在以下各节中直接展开编写。
}\QuickQuizEndE
}

下一步是编写\co{dyntick_irq()}进程来建模中断处理程序：

\input{CodeSamples/formal/promela/dyntick/dyntickRCU-irqnn-ssl@dyntick_irq.fcv}

\begin{fcvref}[ln:formal:promela:dyntick:dyntickRCU-irqnn-ssl:dyntick_irq]
\clnrefrange{do}{od}的循环建模了最多\co{MAX_DYNTICK_LOOP_IRQ}
次中断，\clnref{cond1,cond2}构成循环条件，
\clnref{inc_i}递增控制变量。
\Clnref{in_irq}告诉\co{dyntick_nohz()}中断处理程序正在
运行，\clnref{clr_in_irq}告诉\co{dyntick_nohz()}该处理
程序已完成。
\Clnref{irq_done}用于活性验证，与\co{dyntick_nohz()}中
对应的行相同。
\end{fcvref}

\QuickQuiz{
	\begin{fcvref}[ln:formal:promela:dyntick:dyntickRCU-irqnn-ssl:dyntick_irq]
	为什么\clnref{clr_in_irq,inc_i}（\qtco{in_dyntick_irq = 0;}
	和\qco{i++;}）是原子执行的？
	\end{fcvref}
}\QuickQuizAnswer{
	这些代码行属于控制模型，而不是被建模的代码，
	因此没有理由对它们进行非原子建模。
	将它们原子化建模的动机是减少状态空间的大小。
}\QuickQuizEnd

\begin{fcvref}[ln:formal:promela:dyntick:dyntickRCU-irqnn-ssl:dyntick_irq]
\Clnrefrange{enter:b}{enter:e}建模\co{rcu_irq_enter()}，
\clnref{add_prmt_cnt:b,add_prmt_cnt:e}建模\co{__irq_enter()}
的相关片段。
\Clnrefrange{vrf_safe:b}{vrf_safe:e}以与\co{dynticks_nohz()}
中对应行相同的方式验证安全性。
\Clnref{irq_exit:b,irq_exit:e}建模\co{__irq_exit()}的
相关片段，最后\clnrefrange{exit:b}{exit:e}建模
\co{rcu_irq_exit()}。
\end{fcvref}

\QuickQuiz{
	这个\co{dynticks_irq()}进程无法建模中断的
	哪个属性？
}\QuickQuizAnswer{
	一个这样的属性是嵌套中断，
	下一节将处理这个问题。
}\QuickQuizEnd

\co{grace_period()}进程变为如下所示：

\input{CodeSamples/formal/promela/dyntick/dyntickRCU-irqnn-ssl@grace_period.fcv}

\begin{fcvref}[ln:formal:promela:dyntick:dyntickRCU-irqnn-ssl:grace_period]
\co{grace_period()} 的实现与之前的非常相似。
唯一的变化是在 \clnref{MDLI} 上添加了新的中断计数参数，
在 \clnref{edit1,edit3} 上将新的 \co{dyntick_irq_done}
变量添加到活性检查中，当然还有 \clnref{edit2,edit4} 上的优化。
\end{fcvref}

此模型（\path{dyntickRCU-irqnn-ssl.spin}）
以大约五十万个状态产生了正确的验证，无错误通过。
然而，此版本的模型不处理嵌套中断。
这个主题将在下一节中讨论。

\subsubsection{验证嵌套中断处理程序}
\label{sec:formal:Validating Nested Interrupt Handlers}

嵌套中断处理程序可以通过拆分 \co{dyntick_irq()} 中循环的
主体来建模，如下所示：

\input{CodeSamples/formal/promela/dyntick/dyntickRCU-irq-ssl@dyntick_irq.fcv}

\begin{fcvref}[ln:formal:promela:dyntick:dyntickRCU-irq-ssl:dyntick_irq]
这与之前的 \co{dynticks_irq()} 进程类似。
它在 \clnref{j} 上添加了第二个计数器变量 \co{j}，使得
\co{i} 计数中断处理程序的进入次数，\co{j} 计数退出次数。
\clnref{om} 上的 \co{outermost} 变量帮助确定何时需要
为安全检查采样 \co{grace_period_state} 变量。
\clnref{chk_ex:b,chk_ex:e} 上的循环退出检查被更新为要求
指定数量的中断处理程序既被进入也被退出，
\co{i} 的递增被移到 \clnref{inc_i}，即中断进入模型的末尾。
\Clnrefrange{atm1:b}{atm1:e} 设置 \co{outermost} 变量以指示
这是否是一组嵌套中断中最外层的，
并设置 \co{dyntick_nohz()} 进程使用的
\co{in_dyntick_irq} 变量。
\Clnrefrange{atm2:b}{atm2:e} 捕获 \co{grace_period_state}
变量的状态，但仅在最外层中断处理程序中执行。

\Clnref{cnd_ex} 有中断退出建模的 do 循环条件：
只要我们退出的中断比进入的少，就可以退出另一个中断。
\Clnrefrange{atm3:b}{atm3:e}
检查安全准则，但仅在从最外层中断级别退出时执行。
最后，\clnrefrange{atm4:b}{atm4:e} 递增中断退出计数 \co{j}，
如果这是最外层中断级别，则清除 \co{in_dyntick_irq}。
\end{fcvref}

此模型（\path{dyntickRCU-irq-ssl.spin}）
以略多于五十万个状态产生了正确的验证，无错误通过。
然而，此版本的模型不处理 NMI，
这将在下一节中讨论。

\subsubsection{验证 NMI 处理程序}
\label{sec:formal:Validating NMI Handlers}

我们对 NMI 采取与中断相同的一般方法，
但要记住 NMI 不会嵌套。
这导致了如下所示的 \co{dyntick_nmi()} 进程：

\input{CodeSamples/formal/promela/dyntick/dyntickRCU-irq-nmi-ssl@dyntick_nmi.fcv}

当然，NMI 的存在要求对其他组件进行调整。
例如，\co{EXECUTE_MAINLINE()} 宏现在需要通过检查
\co{dyntick_nmi_done} 变量来同时关注 NMI 处理程序
（\co{in_dyntick_nmi}）和中断处理程序
（\co{in_dyntick_irq}），如下所示：

\input{CodeSamples/formal/promela/dyntick/dyntickRCU-irq-nmi-ssl@EXECUTE_MAINLINE.fcv}

我们还需要引入一个 \co{EXECUTE_IRQ()} 宏，
该宏检查 \co{in_dyntick_nmi} 以允许
\co{dyntick_irq()} 排除 \co{dyntick_nmi()}：

\input{CodeSamples/formal/promela/dyntick/dyntickRCU-irq-nmi-ssl@EXECUTE_IRQ.fcv}

还需要将 \co{dyntick_irq()} 转换为
\co{EXECUTE_IRQ()}，如下所示：

\input{CodeSamples/formal/promela/dyntick/dyntickRCU-irq-nmi-ssl@dyntick_irq.fcv}

\begin{fcvref}[ln:formal:promela:dyntick:dyntickRCU-irq-nmi-ssl:dyntick_irq]
请注意，我们已将 \qco{if} 语句直接展开编写
（例如，\clnrefrange{stmt1:b}{stmt1:e}）。
此外，处理严格局部状态的语句
（如 \clnref{inc_i}）不需要排除 \co{dyntick_nmi()}。
\end{fcvref}

最后，\co{grace_period()} 只需要少量更改：

\input{CodeSamples/formal/promela/dyntick/dyntickRCU-irq-nmi-ssl@grace_period.fcv}

\begin{fcvref}[ln:formal:promela:dyntick:dyntickRCU-irq-nmi-ssl:grace_period]
我们在 \clnref{MDL_NMI} 上为新的 \co{MAX_DYNTICK_LOOP_NMI}
参数添加了 \co{printf()}，并在 \clnref{nmi_done1,nmi_done2}
上将 \co{dyntick_nmi_done} 添加到 \co{shouldexit} 赋值中。
\end{fcvref}

此模型（\path{dyntickRCU-irq-nmi-ssl.spin}）
以数亿个状态产生了正确的验证，无错误通过。

\QuickQuiz{
	Paul \emph{总是}以这种痛苦的增量方式编写代码吗？
}\QuickQuizAnswer{
	不总是，但越来越频繁。
	在这个案例中，Paul 从包含中断处理程序的最小代码片段开始，
	因为他不确定如何最好地在 Promela 中建模中断。
	一旦让它运行起来，他就添加了其他特性。
	（但如果再做一次，他会从一个``玩具''处理程序开始。
	例如，他可能让处理程序将一个变量递增两次，
	并让主线代码验证该值始终为偶数。）

	为什么要用增量方法？
	请考虑以下这段话，出自Brian W. Kernighan：

	\begin{quote}
		调试的难度是编写代码的两倍。
		因此，如果你尽可能巧妙地编写代码，
		那么按定义，你就不够聪明来调试它。
	\end{quote}

	这意味着任何优化代码生产的尝试都应该将至少66\pct\ 的
	重点放在优化调试过程上，即使这会增加编码所花费的时间和
	精力。增量编码和测试是优化调试过程的一种方式，代价是
	编码工作量有所增加。
	Paul使用这种方法是因为他很少有奢侈的条件可以将整天
	（更不用说整周）用于编码和调试。
}\QuickQuizEnd

\subsubsection{（重新）学到的经验}
\label{sec:formal:Lessons (Re)Learned}

这次工作提供了一些（重新）学到的经验教训：

\begin{enumerate}
\item	{\bf Promela 和 Spin 可以验证中断/NMI 处理程序的交互。}
\item	{\bf 为代码编写文档可以帮助定位缺陷。}
	在本例中，文档编写工作定位了
	\co{rcu_enter_nohz()} 和 \co{rcu_exit_nohz()} 中
	一个放错位置的内存屏障，
	如以下补丁所示~\cite{PaulEMcKenney2008commit:ae66be9b71b1}。

\begin{VerbatimU}
 static inline void rcu_enter_nohz(void)
 {
+       mb();
        __get_cpu_var(dynticks_progress_counter)++;
-       mb();
 }

 static inline void rcu_exit_nohz(void)
 {
-       mb();
        __get_cpu_var(dynticks_progress_counter)++;
+       mb();
 }
\end{VerbatimU}

\item	{\bf 尽早、经常、彻底地验证你的代码。}
	这次工作定位了 \co{rcu_try_flip_waitack_needed()} 中
	一个非常微妙的缺陷，该缺陷很难通过测试或调试来发现，
	如以下补丁所示~\cite{PaulEMcKenney2008commit:d7c0651390b6}。

\begin{VerbatimU}
-       if ((curr - snap) > 2 || (snap & 0x1) == 0)
+       if ((curr - snap) > 2 || (curr & 0x1) == 0)
\end{VerbatimU}

\item	{\bf 始终验证你的验证代码。}
	通常的做法是故意插入一个缺陷，
	并验证验证代码能够捕获它。
	当然，如果验证代码未能捕获这个缺陷，
	你可能还需要验证缺陷本身，依此类推，
	无限递归。
	然而，如果你发现自己处于这种境地，
	好好睡一觉可能是一种极其有效的调试技术。
	然后你就会看到，验证验证代码的明显技术是
	在被验证的代码中故意插入缺陷。
	如果验证未能发现它们，那么验证代码显然有问题。
\item	{\bf 使用原子指令可以简化验证。}
	不幸的是，使用 \co{cmpxchg} 原子指令也会减慢
	关键的 \IXacr{irq} 快速路径，因此在本例中不适用。
\item	{\bf 需要复杂的形式验证通常表明需要重新思考你的设计。}
\end{enumerate}

关于最后一点，事实证明 dynticks 问题有一个简单得多的解决方案，
将在下一节中介绍。

\subsubsection{简单性避免了形式验证}
\label{sec:formal:Simplicity Avoids Formal Verification}

可抢占 RCU 的 dynticks 接口的复杂性主要是由于
\IRQ 和 NMI 使用相同的代码路径和相同的状态变量。
这引出了为 \IRQ 和 NMI 提供单独代码路径和变量的想法，
正如分层 RCU~\cite{PaulEMcKenney2008HierarchicalRCU}
所做的那样，这是由
Manfred Spraul~\cite{ManfredSpraul2008StateMachineRCU}
间接建议的。
这项工作在 v2.6.29 开发周期中被合入主线内核~\cite{PaulEMcKenney2008commit:64db4cfff99c}。

\subsubsection{简化 Dynticks 接口的状态变量}
\label{sec:formal:State Variables for Simplified Dynticks Interface}

\Cref{lst:formal:Variables for Simple Dynticks Interface}
显示了新的每 CPU 状态变量。
这些变量被分组到结构体中，以允许多个独立的 RCU 实现
（例如，\co{rcu} 和 \co{rcu_bh}）方便且高效地共享 dynticks 状态。
在以下内容中，它们可以被视为独立的每 CPU 变量。

\begin{listing}
\begin{VerbatimL}
struct rcu_dynticks {
	int dynticks_nesting;
	int dynticks;
	int dynticks_nmi;
};

struct rcu_data {
	...
	int dynticks_snap;
	int dynticks_nmi_snap;
	...
};
\end{VerbatimL}
\caption{简单 Dynticks 接口的变量}
\label{lst:formal:Variables for Simple Dynticks Interface}
\end{listing}

\co{dynticks_nesting}、\co{dynticks} 和 \co{dynticks_snap} 变量
用于 \IRQ\ 代码路径，\co{dynticks_nmi} 和
\co{dynticks_nmi_snap} 变量用于 NMI 代码路径，
尽管 NMI 代码路径也会引用（但不修改）
\co{dynticks_nesting} 变量。
这些变量的使用方式如下：

\begin{description}[style=nextline]
\item	[\tco{dynticks_nesting}]
	这个变量计数对应 CPU 应被监视 RCU 读端临界区的原因数量。
	如果 CPU 处于 dynticks-idle 模式，则它计数
	\IRQ\ 嵌套级别，否则比 \IRQ\ 嵌套级别大一。
\item	[\tco{dynticks}]
	如果对应 CPU 处于 dynticks-idle 模式且该 CPU 上
	没有正在运行的 \IRQ\ 处理程序，则该计数器的值为偶数，
	否则为奇数。
	换句话说，如果该计数器的值为奇数，
	则对应 CPU 可能处于 RCU 读端临界区中。
\item	[\tco{dynticks_nmi}]
	如果对应 CPU 处于 NMI 处理程序中，则该计数器的值为奇数，
	但仅当 NMI 到达时该 CPU 处于 dyntick-idle 模式
	且没有 \IRQ\ 处理程序在运行时才如此。
	否则，该计数器的值将为偶数。
\item	[\tco{dynticks_snap}]
	这将是 \co{dynticks} 计数器的快照，
	但仅当当前 RCU 宽限期已经延续了过长的时间时才如此。
\item	[\tco{dynticks_nmi_snap}]
	这将是 \co{dynticks_nmi} 计数器的快照，
	但同样仅当当前 RCU 宽限期已经延续了过长的时间时才如此。
\end{description}

如果 \co{dynticks} 和 \co{dynticks_nmi} 在给定时间间隔内
都取过偶数值，则对应 CPU 在该间隔内经历了一个\IX{quiescent state}。

\QuickQuiz{
	但如果一个 NMI 处理程序在一个 \IRQ\ 处理程序完成之前
	开始运行，并且该 NMI 处理程序一直运行到第二个
	\IRQ\ 处理程序开始，会怎样？
}\QuickQuizAnswer{
	这在单个 CPU 的范围内不可能发生。
	第一个 \IRQ\ 处理程序在 NMI 处理程序返回之前无法完成。
	因此，如果 \co{dynticks} 和 \co{dynticks_nmi}
	变量在给定时间间隔内都取过偶数值，
	则对应 CPU 确实在该间隔内的某个时刻处于静止状态。
}\QuickQuizEnd

\subsubsection{进入和离开 Dynticks-Idle 模式}
\label{sec:formal:Entering and Leaving Dynticks-Idle Mode}

\Cref{lst:formal:Entering and Exiting Dynticks-Idle Mode}
显示了 \co{rcu_enter_nohz()} 和 \co{rcu_exit_nohz()}，
它们进入和退出 dynticks-idle 模式，也称为``nohz''模式。
这两个函数从进程上下文中调用。

\begin{listing}
\begin{fcvlabel}[ln:formal:Entering and Exiting Dynticks-Idle Mode]
\begin{VerbatimL}[commandchars=\\\[\]]
void rcu_enter_nohz(void)
{
	unsigned long flags;
	struct rcu_dynticks *rdtp;

	smp_mb();			\lnlbl[mb]
	local_irq_save(flags);		\lnlbl[irq_sv]
	rdtp = &__get_cpu_var(rcu_dynticks); \lnlbl[get_ptr]
	rdtp->dynticks++;		\lnlbl[inc_cnt]
	rdtp->dynticks_nesting--;	\lnlbl[dec_nst]
	WARN_ON(rdtp->dynticks & 0x1);
	local_irq_restore(flags);	\lnlbl[irq_rs]
}

void rcu_exit_nohz(void)
{
	unsigned long flags;
	struct rcu_dynticks *rdtp;

	local_irq_save(flags);
	rdtp = &__get_cpu_var(rcu_dynticks);
	rdtp->dynticks++;
	rdtp->dynticks_nesting++;
	WARN_ON(!(rdtp->dynticks & 0x1));
	local_irq_restore(flags);
	smp_mb();
}
\end{VerbatimL}
\end{fcvlabel}
\caption{进入和退出 Dynticks-Idle 模式}
\label{lst:formal:Entering and Exiting Dynticks-Idle Mode}
\end{listing}

\begin{fcvref}[ln:formal:Entering and Exiting Dynticks-Idle Mode]
\Clnref{mb} 确保任何先前的内存访问（可能包括来自 RCU
读端临界区的访问）在标记进入 dynticks-idle 模式之前被其他 CPU 看到。
\Clnref{irq_sv,irq_rs} 禁用和重新启用 \IRQ。
\Clnref{get_ptr} 获取指向当前 CPU 的 \co{rcu_dynticks}
结构的指针，
\clnref{inc_cnt} 递增当前 CPU 的 \co{dynticks} 计数器，
鉴于我们在进程上下文中进入 dynticks-idle 模式，
该计数器现在应该是偶数。
最后，\clnref{dec_nst} 递减 \co{dynticks_nesting}，
现在应该为零。
\end{fcvref}

\co{rcu_exit_nohz()} 函数非常相似，但递增 \co{dynticks_nesting}
而非递减，并检查相反的 \co{dynticks} 极性。

\subsubsection{来自 Dynticks-Idle 模式的 NMI}
\label{sec:formal:NMIs From Dynticks-Idle Mode}

\begin{fcvref}[ln:formal:NMIs From Dynticks-Idle Mode]
\Cref{lst:formal:NMIs From Dynticks-Idle Mode}
显示了 \co{rcu_nmi_enter()} 和 \co{rcu_nmi_exit()} 函数，
它们分别通知 RCU 从 dynticks-idle 模式进入和退出 NMI。
然而，如果 NMI 在 \IRQ\ 处理程序期间到达，
那么 RCU 已经在监视来自该 CPU 的 RCU 读端临界区，
因此如果 \co{dynticks} 为奇数，\co{rcu_nmi_enter()} 的
\clnref{chk_ext1,ret1} 和 \co{rcu_nmi_exit()} 的
\clnref{chk_ext2,ret2} 会静默返回。
否则，两个函数递增 \co{dynticks_nmi}，
\co{rcu_nmi_enter()} 使其保持奇数值，
\co{rcu_nmi_exit()} 使其保持偶数值。
两个函数分别在 \clnref{mb1,mb2} 上在此递增和可能的 RCU
读端临界区之间执行内存屏障。
\end{fcvref}

\begin{listing}
\begin{fcvlabel}[ln:formal:NMIs From Dynticks-Idle Mode]
\begin{VerbatimL}[commandchars=\\\[\]]
void rcu_nmi_enter(void)
{
	struct rcu_dynticks *rdtp;

	rdtp = &__get_cpu_var(rcu_dynticks);
	if (rdtp->dynticks & 0x1)	\lnlbl[chk_ext1]
		return;			\lnlbl[ret1]
	rdtp->dynticks_nmi++;
	WARN_ON(!(rdtp->dynticks_nmi & 0x1));
	smp_mb();			\lnlbl[mb1]
}

void rcu_nmi_exit(void)
{
	struct rcu_dynticks *rdtp;

	rdtp = &__get_cpu_var(rcu_dynticks);
	if (rdtp->dynticks & 0x1)	\lnlbl[chk_ext2]
		return;			\lnlbl[ret2]
	smp_mb();			\lnlbl[mb2]
	rdtp->dynticks_nmi++;
	WARN_ON(rdtp->dynticks_nmi & 0x1);
}
\end{VerbatimL}
\end{fcvlabel}
\caption{来自 Dynticks-Idle 模式的 NMI}
\label{lst:formal:NMIs From Dynticks-Idle Mode}
\end{listing}

\subsubsection{来自 Dynticks-Idle 模式的中断}
\label{sec:formal:Interrupts From Dynticks-Idle Mode}

\begin{fcvref}[ln:formal:Interrupts From Dynticks-Idle Mode]
\Cref{lst:formal:Interrupts From Dynticks-Idle Mode}
显示了 \co{rcu_irq_enter()} 和 \co{rcu_irq_exit()}，
它们分别通知 RCU 进入和退出 \IRQ\ 上下文。
\co{rcu_irq_enter()} 的 \clnref{inc_nst} 递增 \co{dynticks_nesting}，
如果该变量已经非零，\clnref{ret1} 会静默返回。
否则，\clnref{inc_dynt1} 递增 \co{dynticks}，
该计数器随后将具有奇数值，与该 CPU 现在可以执行
RCU 读端临界区的事实一致。
因此 \clnref{mb1} 执行一个内存屏障以确保
\co{dynticks} 的递增在后续 \IRQ\ 处理程序可能执行的
任何 RCU 读端临界区之前被看到。

\begin{listing}
\begin{fcvlabel}[ln:formal:Interrupts From Dynticks-Idle Mode]
\begin{VerbatimL}[commandchars=\\\[\]]
void rcu_irq_enter(void)
{
	struct rcu_dynticks *rdtp;

	rdtp = &__get_cpu_var(rcu_dynticks);
	if (rdtp->dynticks_nesting++)		\lnlbl[inc_nst]
		return;				\lnlbl[ret1]
	rdtp->dynticks++;			\lnlbl[inc_dynt1]
	WARN_ON(!(rdtp->dynticks & 0x1));
	smp_mb();				\lnlbl[mb1]
}

void rcu_irq_exit(void)
{
	struct rcu_dynticks *rdtp;

	rdtp = &__get_cpu_var(rcu_dynticks);
	if (--rdtp->dynticks_nesting)		\lnlbl[dec_nst]
		return;				\lnlbl[ret2]
	smp_mb();				\lnlbl[mb2]
	rdtp->dynticks++;			\lnlbl[inc_dynt2]
	WARN_ON(rdtp->dynticks & 0x1);		\lnlbl[chk_even]
	if (__get_cpu_var(rcu_data).nxtlist ||	\lnlbl[chk_cb:b]
	    __get_cpu_var(rcu_bh_data).nxtlist)
		set_need_resched();		\lnlbl[chk_cb:e]
}
\end{VerbatimL}
\end{fcvlabel}
\caption{来自 Dynticks-Idle 模式的中断}
\label{lst:formal:Interrupts From Dynticks-Idle Mode}
\end{listing}

\co{rcu_irq_exit()} 的 \clnref{dec_nst} 递减 \co{dynticks_nesting}，
如果结果非零，\clnref{ret2} 静默返回。
否则，\clnref{mb2} 执行内存屏障以确保
\clnref{inc_dynt2} 上 \co{dynticks} 的递增在先前
\IRQ\ 处理程序可能执行的任何 RCU 读端临界区之后被看到。
\Clnref{chk_even} 验证 \co{dynticks} 现在为偶数，
与 dynticks-idle 模式中不能出现 RCU 读端临界区的事实一致。
\Clnrefrange{chk_cb:b}{chk_cb:e} 检查先前的 \IRQ\ 处理程序
是否排入了任何 RCU 回调，如果是则通过重调度 API
强制此 CPU 退出 dynticks-idle 模式。
\end{fcvref}

\subsubsection{检查 Dynticks 静止状态}
\label{sec:formal:Checking For Dynticks Quiescent States}

\begin{fcvref}[ln:formal:Saving Dyntick Progress Counters]
\Cref{lst:formal:Saving Dyntick Progress Counters}
显示了 \co{dyntick_save_progress_counter()}，它获取
指定 CPU 的 \co{dynticks} 和 \co{dynticks_nmi} 计数器的快照。
\Clnref{snap,snapn} 将这两个变量快照到局部变量中，
\clnref{mb} 执行内存屏障以与
\cref{lst:formal:Entering and Exiting Dynticks-Idle Mode,%
lst:formal:NMIs From Dynticks-Idle Mode,%
lst:formal:Interrupts From Dynticks-Idle Mode}
中函数的内存屏障配对。
\Clnref{rec_snap,rec_snapn} 记录快照以供后续调用
\co{rcu_implicit_dynticks_qs()} 使用，
\clnref{chk_prgs} 检查 CPU 是否处于 dynticks-idle 模式且
没有正在进行的 \IRQ 或 NMI（换句话说，两个快照都是偶数值），
因此处于扩展静止状态。
如果是，\clnref{cnt:b,cnt:e} 计数此事件，
\clnref{ret} 在 CPU 处于静止状态时返回 true。
\end{fcvref}

\begin{listing}
\begin{fcvlabel}[ln:formal:Saving Dyntick Progress Counters]
\begin{VerbatimL}[commandchars=\\\[\]]
static int
dyntick_save_progress_counter(struct rcu_data *rdp)
{
	int ret;
	int snap;
	int snap_nmi;

	snap = rdp->dynticks->dynticks;		\lnlbl[snap]
	snap_nmi = rdp->dynticks->dynticks_nmi;	\lnlbl[snapn]
	smp_mb();				\lnlbl[mb]
	rdp->dynticks_snap = snap;		\lnlbl[rec_snap]
	rdp->dynticks_nmi_snap = snap_nmi;	\lnlbl[rec_snapn]
	ret = ((snap & 0x1) == 0) && ((snap_nmi & 0x1) == 0); \lnlbl[chk_prgs]
	if (ret)				\lnlbl[cnt:b]
		rdp->dynticks_fqs++;		\lnlbl[cnt:e]
	return ret;				\lnlbl[ret]
}
\end{VerbatimL}
\end{fcvlabel}
\caption{保存 Dyntick 进度计数器}
\label{lst:formal:Saving Dyntick Progress Counters}
\end{listing}

\begin{fcvref}[ln:formal:Checking Dyntick Progress Counters]
\Cref{lst:formal:Checking Dyntick Progress Counters}
显示了 \co{rcu_implicit_dynticks_qs()}，它被调用以检查
CPU 是否在调用 \co{dynticks_save_progress_counter()} 之后
进入了 dyntick-idle 模式。
\Clnref{curr,currn} 获取对应 CPU 的 \co{dynticks} 和
\co{dynticks_nmi} 变量的新快照，
而 \clnref{snap,snapn} 检索之前由
\co{dynticks_save_progress_counter()} 保存的快照。
\Clnref{mb} 然后执行内存屏障以与
\cref{lst:formal:Entering and Exiting Dynticks-Idle Mode,%
lst:formal:NMIs From Dynticks-Idle Mode,%
lst:formal:Interrupts From Dynticks-Idle Mode}
中函数的内存屏障配对。
\Clnrefrange{chk_q:b}{chk_q:e} 然后检查 CPU 是否当前
处于静止状态（\co{curr} 和 \co{curr_nmi} 具有偶数值），
或者自上次调用 \co{dynticks_save_progress_counter()}
以来是否经历了静止状态（\co{dynticks} 和 \co{dynticks_nmi}
的值已改变）。
如果这些检查确认 CPU 已经经历了 dyntick-idle 静止状态，
则 \clnref{cnt} 计数该事实，
\clnref{ret_1} 返回该事实的指示。
无论如何，\clnref{chk_race} 检查可能导致 RCU
等待已离线 CPU 的竞态条件。
\end{fcvref}

\begin{listing}
\begin{fcvlabel}[ln:formal:Checking Dyntick Progress Counters]
\begin{VerbatimL}[commandchars=\\\[\]]
static int
rcu_implicit_dynticks_qs(struct rcu_data *rdp)
{
	long curr;
	long curr_nmi;
	long snap;
	long snap_nmi;

	curr = rdp->dynticks->dynticks;		\lnlbl[curr]
	snap = rdp->dynticks_snap;		\lnlbl[snap]
	curr_nmi = rdp->dynticks->dynticks_nmi;	\lnlbl[currn]
	snap_nmi = rdp->dynticks_nmi_snap;	\lnlbl[snapn]
	smp_mb();				\lnlbl[mb]
	if ((curr != snap || (curr & 0x1) == 0) && \lnlbl[chk_q:b]
	    (curr_nmi != snap_nmi || (curr_nmi & 0x1) == 0)) { \lnlbl[chk_q:e]
		rdp->dynticks_fqs++;		\lnlbl[cnt]
		return 1;			\lnlbl[ret_1]
	}
	return rcu_implicit_offline_qs(rdp);	\lnlbl[chk_race]
}
\end{VerbatimL}
\end{fcvlabel}
\caption{检查 Dyntick 进度计数器}
\label{lst:formal:Checking Dyntick Progress Counters}
\end{listing}

\QuickQuiz{
	这仍然相当复杂。
	为什么不直接使用一个带有每 CPU 位的 \co{cpumask_t}，
	在进入 \IRQ\ 或 NMI 处理程序时清除该位，
	在退出时设置它？
}\QuickQuizAnswer{
	虽然这种方法在功能上是正确的，但它会导致
	在大型机器上产生过多的 \IRQ\ 进入/退出开销。
	相比之下，本节中阐述的方法允许每个 CPU 在
	\IRQ\ 和 NMI 进入/退出时只触碰每 CPU 数据，
	从而产生更低的 \IRQ\ 进入/退出开销，
	尤其是在大型机器上。
}\QuickQuizEnd

Linux 内核 RCU 的 dyntick-idle 代码后来又基于
Andy Lutomirski~\cite{PaulMcKenney2015dyntickAndyNMI}
的建议被重写了一次，
但现在是时候总结并继续讨论其他主题了。

\subsubsection{讨论}
\label{sec:formal:Discussion}

一个小小的观点转变导致了 RCU dynticks 接口的实质性简化。
导致这种简化的关键变化是将 \IRQ\ 和 NMI 上下文之间的共享减到最小。
在这个简化接口中，唯一的共享是 NMI 上下文对
\IRQ\ 变量（\co{dynticks} 变量）的引用。
这种类型的共享是良性的，因为 NMI 函数从不更新该变量，
因此其值在 NMI 处理程序的整个生命周期中保持不变。
这种共享的限制使得各个函数可以逐个理解，
与\cref{sec:formal:Promela Parable: dynticks and Preemptible RCU}
中描述的情况形成鲜明对比——在那里，NMI 可能在
\IRQ\ 函数执行期间的任何时刻修改共享状态。

验证是一件非常好的事，但简单性更好。

% Restore frame around VerbatimN in two-column layout
\IfTwoColumn{
\RecustomVerbatimEnvironment{VerbatimN}{Verbatim}%
{numbers=left,numbersep=3pt,xleftmargin=5pt,xrightmargin=5pt,frame=single}
\setlength{\lnlblraise}{6pt}
}{}

% formal/ppcmem.tex
% mainfile: ../perfbook.tex
% SPDX-License-Identifier: CC-BY-SA-3.0

\section{专用状态空间搜索}
\label{sec:formal:Special-Purpose State-Space Search}
%
\epigraph{样样通，样样松。}{佚名}

虽然 Promela 和 Spin 允许你验证几乎任何（较小的）算法，
但它们的通用性有时会成为一种负担。
例如，Promela 不理解 \IX{memory model} 或任何类型的
重排序语义。
因此，本节描述了一些理解生产系统使用的 memory model 的
state-space search 工具，大大简化了弱序代码的验证。

例如，
\cref{sec:formal:Promela Example: QRCU}
展示了如何说服 Promela 来考虑弱内存排序。
虽然这种方法可以很好地工作，但它要求开发者
完全理解系统的 memory model。
不幸的是，很少有（如果有的话）开发者完全理解
现代 CPU 的复杂 memory model。

因此，另一种方法是使用已经理解内存排序的工具，
例如由剑桥大学的 \ppl{Peter}{Sewell} 和 \ppl{Susmit}{Sarkar}、
INRIA 的 \ppl{Luc}{Maranget}、\ppl{Francesco Zappa}{Nardelli}
和 \ppl{Pankaj}{Pawan}、牛津大学的 \ppl{Jade}{Alglave}
与 IBM 的 \ppl{Derek}{Williams} 合作产生的
PPCMEM 工具~\cite{JadeAlglave2011ppcmem}。
这个团队形式化了 Power、\ARM、x86 以及
C/C++11 标准~\cite{RichardSmith2019N4800}的 memory model，
并基于 Power 和 \ARM\ 形式化产生了 PPCMEM 工具。

\QuickQuiz{
	但 x86 有强内存排序，为什么要形式化它的
	memory model？
}\QuickQuizAnswer{
	实际上，学术界认为 x86 的 memory model 是弱的，
	因为它可以允许先前的存储与后续的加载重排序。
	从学术角度来看，强 memory model 是完全不允许
	重排序的，这样所有线程都同意它们可见的所有操作的顺序。

	而且确实存在开发者有时对 x86 内存排序
	感到困惑的情况。
}\QuickQuizEnd

PPCMEM 工具以 \emph{litmus test} 作为输入。
一个示例 litmus test 在
\cref{sec:formal:Anatomy of a Litmus Test}中介绍。
\Cref{sec:formal:What Does This Litmus Test Mean?}
将这个 litmus test 与等效的 C 语言程序联系起来，
\cref{sec:formal:Running a Litmus Test}描述了如何
将 PPCMEM 应用于这个 litmus test，而
\cref{sec:formal:PPCMEM Discussion}
讨论了其影响。

\subsection{Litmus Test 的剖析}
\label{sec:formal:Anatomy of a Litmus Test}

\cref{lst:formal:PPCMEM Litmus Test}展示了
PPCMEM 的一个 PowerPC litmus test 示例。
ARM 接口以相同方式工作，但用 \ARM\ 指令
替代 Power 指令，并将开头的 \qco{PPC}
替换为 \qco{ARM}。

\begin{listing}
\begin{fcvlabel}[ln:formal:PPCMEM Litmus Test]
\begin{VerbatimL}[commandchars=\@\[\]]
PPC SB+lwsync-RMW-lwsync+isync-simple		@lnlbl[type]
""						@lnlbl[altname]
{						@lnlbl[init:b]
0:r2=x; 0:r3=2; 0:r4=y; 0:r10=0; 0:r11=0; 0:r12=z; @lnlbl[init:0]
1:r2=y; 1:r4=x;					@lnlbl[init:1]
}						@lnlbl[init:e]
 P0                 | P1           ;		@lnlbl[procid]
 li r1,1            | li r1,1      ;		@lnlbl[reginit]
 stw r1,0(r2)       | stw r1,0(r2) ;		@lnlbl[stw]
 lwsync             | sync         ; @lnlbl[P0lwsync] @lnlbl[P1sync]
                    | lwz r3,0(r4) ; @lnlbl[P0empty]  @lnlbl[P1lwz]
 lwarx  r11,r10,r12 | ;		@lnlbl[P0lwarx] @lnlbl[P1empty:b]
 stwcx. r11,r10,r12 | ;		@lnlbl[P0stwcx]
 bne Fail1          | ;		@lnlbl[P0bne]
 isync              | ;		@lnlbl[P0isync]
 lwz r3,0(r4)       | ;		@lnlbl[P0lwz]
 Fail1:             | ;		@lnlbl[P0fail1] @lnlbl[P1empty:e]

exists						@lnlbl[assert:b]
(0:r3=0 /\ 1:r3=0)				@lnlbl[assert:e]
\end{VerbatimL}
\end{fcvlabel}
\caption{PPCMEM Litmus Test}
\label{lst:formal:PPCMEM Litmus Test}
\end{listing}

\begin{fcvref}[ln:formal:PPCMEM Litmus Test]
在这个例子中，\clnref{type}标识了系统类型（\qco{ARM} 或
\qco{PPC}）并包含模型的标题。
\Clnref{altname}提供了一个测试的替代名称位置，
通常你会像上面的例子那样留空。
可以使用 OCaml（或 Pascal）的 \nbco{(* *)} 语法
在\clnref{altname,init:b}之间插入注释。

\Clnrefrange{init:b}{init:e}给出所有寄存器的初始值；
每个的形式为
\co{P:R=V}，其中 \co{P} 是进程标识符，\co{R} 是寄存器
标识符，\co{V} 是值。
例如，进程~0 的寄存器 \co{r3} 初始包含值~2。
如果值是一个变量（在例子中是 \co{x}、\co{y} 或 \co{z}），
那么寄存器被初始化为该变量的地址。
也可以初始化变量的内容，例如，
\co{x=1} 将 \co{x} 的值初始化为~1。
未初始化的变量默认为零值，所以在
例子中，\co{x}、\co{y} 和 \co{z} 初始都为零。

\Clnref{procid}提供了两个进程的标识符，所以
\clnref{init:0}上的 \co{0:r3=2} 也可以写成
\co{P0:r3=2}。
\Clnref{procid}是必须的，标识符必须是
\co{Pn} 的形式，其中 \co{n} 是列号，从最左列的零开始。
这可能看起来不必要地严格，但它确实防止了
实际使用中相当多的混淆。
\end{fcvref}

\QuickQuiz{
	\begin{fcvref}[ln:formal:PPCMEM Litmus Test]
	为什么\cref{lst:formal:PPCMEM Litmus Test}的
	\clnref{reginit}要初始化寄存器？
	为什么不在\clnref{init:0,init:1}上初始化它们？
	\end{fcvref}
}\QuickQuizAnswer{
	两种方式都可以。
	然而，一般来说，使用初始化比使用显式指令更好。
	在这个例子中使用显式指令是为了演示它们的使用。
	此外，该工具网站上提供的许多 litmus test
	（\url{https://www.cl.cam.ac.uk/~pes20/ppcmem/}）
	是自动生成的，这会生成显式的初始化指令。
}\QuickQuizEnd

\begin{fcvref}[ln:formal:PPCMEM Litmus Test]
\Clnrefrange{reginit}{P0fail1}是每个进程的代码行。
一个给定的进程可以有空行，如 P0 的
\clnref{P0empty}和 P1 的\clnrefrange{P1empty:b}{P1empty:e}。
标签和分支是允许的，如\clnref{P0bne}上的分支
到\clnref{P0fail1}上的标签所演示的那样。
话虽如此，过于自由地使用分支会扩展 state space。
使用循环是一种特别好的方式来让你的 state space 爆炸。

\Clnrefrange{assert:b}{assert:e}展示了 assertion，在这种情况下
表示我们关心的是在两个线程完成执行后，P0 和 P1 的 \co{r3} 寄存器
是否都能包含零。
这个 assertion 很重要，因为有许多用例
如果 P0 和 P1 在各自的 \co{r3} 寄存器中都看到零，
将会惨败。

这应该给你足够的信息来构造简单的 litmus test。
还有一些额外的文档可用，不过大部分
额外文档是为另一个在实际硬件上运行测试的
研究工具准备的。
也许更重要的是，在线工具提供了大量预先存在的 litmus test
（可通过 \url{https://www.cl.cam.ac.uk/~pes20/ppcmem/}
的 ``Select ARM Test'' 和 ``Select POWER Test'' 按钮获得）。
这些预先存在的 litmus test 中很可能有一个能
回答你的 Power 或 \ARM\ 内存排序问题。

\subsection{这个 Litmus Test 是什么意思？}
\label{sec:formal:What Does This Litmus Test Mean?}

P0 的\clnref{reginit,stw}等价于 C 语句 \co{x=1}，
因为\clnref{init:0}将 P0 的寄存器 \co{r2} 定义为
\co{x} 的地址。
P0 的\clnref{P0lwarx,P0stwcx}分别是 load-linked
（在 \ARM\ 中称为 ``load register exclusive''，
在 Power 中称为 ``load reserve''）
和 store-conditional（在 \ARM\ 中称为
``store register exclusive''）的助记符。
当它们一起使用时，形成一个原子指令序列，
大致类似于以 x86 \co{lock;cmpxchg} 指令为代表的
\IXacrml{cas} 序列。
提升到更高的抽象层次，从
\clnrefrange{P0lwsync}{P0isync}的序列
等价于 Linux 内核的 \co{atomic_add_return(&z, 0)}。
最后，\clnref{P0lwz}大致等价于 C 语句 \co{r3=y}。

P1 的\clnref{reginit,stw}等价于 C 语句 \co{y=1}，
\clnref{P1sync}是一个内存屏障，等价于 Linux 内核语句 \co{smp_mb()}，
\clnref{P1lwz}等价于 C 语句 \co{r3=x}。
\end{fcvref}

\QuickQuiz{
	\begin{fcvref}[ln:formal:PPCMEM Litmus Test]
	但是\cref{lst:formal:PPCMEM Litmus Test}的
	\clnref{P0fail1}，即 \co{Fail1:} 标签那行，
	怎么了？
	\end{fcvref}
}\QuickQuizAnswer{
	PowerPC 版本的 \co{atomic_add_return()} 实现
	在 \co{stwcx} 指令失败时循环，它通过在
	条件码寄存器中设置非零状态来通知这一点，
	然后由 \co{bne} 指令测试。
	因为实际建模循环会导致 state-space 爆炸，
	我们改为跳转到 \co{Fail1:} 标签，
	以 P0 的 \co{r3} 寄存器中的初始值~2 终止模型，
	这不会触发 exists assertion。

	关于这个技巧是否普遍适用存在一些争论，
	但我还没有看到它失败的例子。
}\QuickQuizEnd

\begin{listing}
\begin{VerbatimL}
void P0(void)
{
	int r3;

	x = 1; /* Lines 8 and 9 */
	atomic_add_return(&z, 0); /* Lines 10-15 */
	r3 = y; /* Line 16 */
}

void P1(void)
{
	int r3;

	y = 1; /* Lines 8-9 */
	smp_mb(); /* Line 10 */
	r3 = x; /* Line 11 */
}
\end{VerbatimL}
\caption{PPCMEM Litmus Test 的含义}
\label{lst:formal:Meaning of PPCMEM Litmus Test}
\end{listing}

将所有这些放在一起，整个 litmus test 的 C 语言等价物
如\cref{lst:formal:Meaning of PPCMEM Litmus Test}所示。
关键点是如果 \co{atomic_add_return()} 充当完整的
内存屏障（如 Linux 内核要求的那样），
那么 \co{P0()} 和 \co{P1()} 的 \co{r3}
变量在执行完成后应该不可能都为零。

下一节描述如何运行这个 litmus test。

\subsection{运行 Litmus Test}
\label{sec:formal:Running a Litmus Test}

如前所述，litmus test 可以通过
\url{https://www.cl.cam.ac.uk/~pes20/ppcmem/}交互式运行，
这有助于建立对 \IX{memory model} 的理解。
然而，这种方法要求用户手动进行
完整的 state-space search。
因为很难确保你已经检查了每种可能的事件序列，
所以为此目的提供了一个单独的工具~\cite{PaulEMcKenney2011ppcmem}。

\begin{listing}
\begin{VerbatimL}[numbers=none,xleftmargin=0pt]
./ppcmem -model lwsync_read_block \
         -model coherence_points filename.litmus
...
States 6
0:r3=0; 1:r3=0;
0:r3=0; 1:r3=1;
0:r3=1; 1:r3=0;
0:r3=1; 1:r3=1;
0:r3=2; 1:r3=0;
0:r3=2; 1:r3=1;
Ok
Condition exists (0:r3=0 /\ 1:r3=0)
Hash=e2240ce2072a2610c034ccd4fc964e77
Observation SB+lwsync-RMW-lwsync+isync Sometimes 1
\end{VerbatimL}
\caption{PPCMEM 检测到错误}
\label{lst:formal:PPCMEM Detects an Error}
\end{listing}

因为\cref{lst:formal:PPCMEM Litmus Test}中
展示的 litmus test 包含读-修改-写指令，我们必须在命令行
中添加 \co{-model} 参数。
如果 litmus test 存储在 \co{filename.litmus} 中，
结果将产生如\cref{lst:formal:PPCMEM Detects an Error}所示的输出，
其中 \co{...} 代表大量的进度输出。
状态列表包括 \co{0:r3=0; 1:r3=0;}，再次表明
旧的 PowerPC 实现的 \co{atomic_add_return()} 不能
充当完整屏障。
最后一行的 ``Sometimes'' 确认了这一点：
assertion 在某些执行中触发，但不是每次都触发。

\begin{listing}
\begin{VerbatimL}[numbers=none,xleftmargin=0pt]
./ppcmem -model lwsync_read_block \
         -model coherence_points filename.litmus
...
States 5
0:r3=0; 1:r3=1;
0:r3=1; 1:r3=0;
0:r3=1; 1:r3=1;
0:r3=2; 1:r3=0;
0:r3=2; 1:r3=1;
No (allowed not found)
Condition exists (0:r3=0 /\ 1:r3=0)
Hash=77dd723cda9981248ea4459fcdf6097d
Observation SB+lwsync-RMW-lwsync+sync Never 0 5
\end{VerbatimL}
\caption{修复后的 Litmus Test 上的 PPCMEM}
\label{lst:formal:PPCMEM on Repaired Litmus Test}
\end{listing}

修复这个 Linux 内核 bug 的方法是将 P0 的 \co{isync} 替换为
\co{sync}，这产生了如
\cref{lst:formal:PPCMEM on Repaired Litmus Test}所示的输出。
如你所见，\co{0:r3=0; 1:r3=0;} 没有出现在状态列表中，
最后一行标注了 ``Never''。
因此，模型预测有问题的执行序列不会发生。

\QuickQuizSeries{%
\QuickQuizB{
	\ARM\ Linux 内核有类似的 bug 吗？
}\QuickQuizAnswer{
	\ARM\ 没有这个特定的 bug，因为它在
	\co{atomic_add_return()} 函数的汇编语言实现
	前后都放置了 \co{smp_mb()}。
	PowerPC 也不再有这个 bug；它早已被
	修复~\cite{BenjaminHerrenschmidt2011:powerpc:atomic_return}。
}\QuickQuizEndB
%
\QuickQuizE{
	\begin{fcvref}[ln:formal:PPCMEM Litmus Test]
	\cref{lst:formal:PPCMEM Litmus Test}中
	\clnref{P0lwsync}上的 \co{lwsync} 是否提供了
	足够的排序？
	\end{fcvref}
}\QuickQuizAnswerE{
	这取决于所需的语义。
	本回答的其余部分假设
	\cref{lst:formal:PPCMEM Litmus Test}中
	\co{P0} 的汇编语言旨在实现一个返回值的原子操作。

	正如\cref{chp:Advanced Synchronization: Memory Ordering}中
	所讨论的，
	Linux 内核的内存一致性模型要求
	返回值的原子 RMW 操作在两侧都是完全排序的。
	\co{lwsync} 提供的排序不足以达到此目的，
	因此应该使用 \co{sync} 代替。
	这个更改后来已经
	完成~\cite{BoqunFeng2015:powerpc:value-returning-atomics}，
	以响应讨论其他几个 litmus test 的
	邮件线程~\cite{Paulmck2015:powerpc:value-returning-atomics}。
	找到 Linux 内核可能存在的其他 bug 留作
	读者的练习。

	在提供较弱语义的其他环境中，\co{lwsync}
	可能就够了。
	但对于 Linux 内核的返回值原子操作则不行！
}\QuickQuizEndE
}

\subsection{PPCMEM 讨论}
\label{sec:formal:PPCMEM Discussion}

这些工具有望对从事在 \ARM\ 和 Power 上运行的底层
并行原语工作的人员提供巨大帮助。
这些工具确实有一些固有的局限性：

\begin{enumerate}
\item	这些工具是研究原型，因此不提供支持。
\item	这些工具不构成 IBM 或 \ARM\ 对其各自 CPU 架构的
	官方声明。
	例如，两家公司保留在任何时间对这些工具的任何版本
	报告 bug 的权利。
	因此，这些工具不能替代在真实硬件上进行仔细的
	压力测试。
	此外，工具和它们所基于的模型都在积极开发中，
	可能随时发生变化。
	另一方面，这个模型是与相关硬件专家协商开发的，
	所以有很好的理由相信它是对架构的可靠表示。
\item	这些工具目前处理指令集的一个子集。
	这个子集对我的目的来说已经足够了，但你的情况
	可能不同。
	特别是，该工具只处理字大小的访问（32位），
	并且被访问的字必须正确对齐。\footnote{
		但最近的工作关注混合大小的
		访问~\cite{Flur:2017:MCA:3093333.3009839}。}
	此外，该工具不处理 \ARM\ 内存屏障指令的某些
	较弱变体，也不处理算术运算。
\item	这些工具仅限于在少量线程上运行的小型无循环代码片段。
	更大的例子会导致 state-space 爆炸，
	就像 Promela 和 Spin 等类似工具一样。
\item	完整的 state-space search 不会指出
	每个有问题的状态是如何达到的。
	话虽如此，一旦你意识到该状态实际上是可达的，
	使用交互式工具找到该状态通常不会太难。
\item	这些工具对复杂的数据结构不太擅长，虽然
	可以使用形如 \qco{x=y; y=z; z=42;} 的
	初始化语句创建和遍历极其简单的链表。
\item	这些工具不处理内存映射 I/O 或设备寄存器。
	当然，处理这些东西需要将它们形式化，
	而这似乎不太可能。
\item	这些工具只会检测你为之编写 assertion 的问题。
	这个弱点是所有形式化方法所共有的，也是测试
	仍然重要的又一个原因。
	用本章开头引用的 Donald Knuth 的不朽名言来说，
	``当心上述代码中的 bug；
	我只是证明了它是正确的，而没有试过。''
\end{enumerate}

话虽如此，这些工具的一个优势是它们被设计为
模拟架构所允许的全部行为范围，包括
合法但当前硬件实现尚未施加于不知情的软件
开发者身上的行为。
因此，经过这些工具审查的算法在实际硬件上运行时
可能具有一些额外的安全边际。
此外，在真实硬件上的测试只能找到 bug；这样的测试
本质上无法证明给定的用法是正确的。
要理解这一点，考虑到研究人员在真实硬件上
常规运行超过1000亿次测试来验证他们的模型。
在一个案例中，尽管运行了1760亿次，
架构允许的行为没有出现~\cite{JadeAlglave2011ppcmem}。
相比之下，完整的 state-space search 允许工具
证明代码片段是正确的。

值得重复的是，形式化方法和工具不能替代测试。
事实是，生产大型可靠的并发软件工件（例如 Linux 内核）
是相当困难的。
因此，开发者必须准备好运用他们所掌握的每一个工具
来实现这个目标。
本章介绍的工具能够定位通过测试非常难以产生
（更不用说追踪）的 bug。
另一方面，测试可以应用于比本章介绍的工具
可能处理的大得多的软件体。
一如既往，使用合适的工具做合适的工作！

当然，最好是通过设计你的并行代码使其容易分区，
然后使用更高级别的原语（如锁、序列计数器、
原子操作和 RCU）来更直接地完成工作，
从而避免在这个层面工作的需要。
即使你绝对必须使用底层的内存屏障和读-修改-写指令
来完成工作，你越保守地使用这些锋利的工具，
你的生活就越可能轻松。

% formal/axiomatic.tex
% mainfile: ../perfbook.tex
% SPDX-License-Identifier: CC-BY-SA-3.0

\section{公理化方法}
\label{sec:formal:Axiomatic Approaches}
\OriginallyPublished{sec:formal:Axiomatic Approaches}{Axiomatic Approaches}{Linux Weekly News}{PaulEMcKenney2014weakaxiom}
%
\epigraph{理论帮助我们承受对事实的无知。}
	{George Santayana}

虽然 PPCMEM 工具可以求解如
\cref{lst:formal:IRIW Litmus Test}所示的著名
``independent reads of independent writes''（IRIW）litmus test，
但这样做需要不少于14个 CPU 小时，
并产生不少于10 GB 的 state space。
话虽如此，这种情况相比 PPCMEM 出现之前已经大有改善，
在那之前解决这个问题需要翻阅大量参考手册、
尝试证明、与专家讨论，
然后在所有这些之后，仍然怀疑最终答案。
虽然14小时看起来可能很长，但它比几周甚至几个月
要短得多。

\begin{listing}
\begin{fcvlabel}[ln:formal:IRIW Litmus Test]
\begin{VerbatimL}[commandchars=\%\@\$]
PPC IRIW.litmus
""
(* Traditional IRIW. *)
{
0:r1=1; 0:r2=x;
1:r1=1;         1:r4=y;
        2:r2=x; 2:r4=y;
        3:r2=x; 3:r4=y;
}
P0           | P1           | P2           | P3           ;
stw r1,0(r2) | stw r1,0(r4) | lwz r3,0(r2) | lwz r3,0(r4) ;
             |              | sync         | sync         ;
             |              | lwz r5,0(r4) | lwz r5,0(r2) ;

exists
(2:r3=1 /\ 2:r5=0 /\ 3:r3=1 /\ 3:r5=0)
\end{VerbatimL}
\end{fcvlabel}
\caption{IRIW Litmus Test}
\label{lst:formal:IRIW Litmus Test}
\end{listing}

然而，鉴于 litmus test 的简单性，所需的时间和空间
有点令人惊讶——它有两个线程存储到两个不同的变量，
另外两个线程以相反的顺序从这两个变量加载。
如果两个加载线程对两个存储的顺序有分歧，
则 assertion 触发。
即使按照 memory-order litmus test 的标准，这也是相当简单的。

PPCMEM 消耗大量时间和空间的一个原因是它进行了
基于轨迹的完整 state-space search，这意味着它必须生成
和评估在架构级别上事件和动作的所有可能顺序和组合。
在这个级别上，加载和存储都对应于精心设计的
事件和动作序列，导致必须完全搜索的非常大的 state space，
进而导致大量的内存和 CPU 消耗。

当然，许多轨迹彼此非常相似，这表明
将相似轨迹视为一个的方法可能会
提高性能。
一种这样的方法是 \pplsur{Jade}{Alglave}
等人~\cite{Alglave:2014:HCM:2594291.2594347}
的 axiomatic approach，它创建一组公理来表示 \IX{memory model}，
然后将 litmus test 转换为可以在这组公理上
证明或反证的定理。
由此产生的工具称为 \qco{herd}，方便地接受与 PPCMEM
相同的 litmus test 作为输入，包括
\cref{lst:formal:IRIW Litmus Test}中所示的 IRIW litmus test。

\begin{listing}
\begin{fcvlabel}[ln:formal:Expanded IRIW Litmus Test]
\begin{VerbatimL}[commandchars=\%\@\$]
PPC IRIW5.litmus
""
(* Traditional IRIW, but with five stores instead of *)
(* just one.                                         *)
{
0:r1=1; 0:r2=x;
1:r1=1;         1:r4=y;
        2:r2=x; 2:r4=y;
        3:r2=x; 3:r4=y;
}
P0           | P1           | P2           | P3           ;
stw r1,0(r2) | stw r1,0(r4) | lwz r3,0(r2) | lwz r3,0(r4) ;
addi r1,r1,1 | addi r1,r1,1 | sync         | sync         ;
stw r1,0(r2) | stw r1,0(r4) | lwz r5,0(r4) | lwz r5,0(r2) ;
addi r1,r1,1 | addi r1,r1,1 |              |              ;
stw r1,0(r2) | stw r1,0(r4) |              |              ;
addi r1,r1,1 | addi r1,r1,1 |              |              ;
stw r1,0(r2) | stw r1,0(r4) |              |              ;
addi r1,r1,1 | addi r1,r1,1 |              |              ;
stw r1,0(r2) | stw r1,0(r4) |              |              ;

exists
(2:r3=1 /\ 2:r5=0 /\ 3:r3=1 /\ 3:r5=0)
\end{VerbatimL}
\end{fcvlabel}
\caption{扩展的 IRIW Litmus Test}
\label{lst:formal:Expanded IRIW Litmus Test}
\end{listing}

然而，PPCMEM 需要14个 CPU 小时来求解 IRIW，\co{herd} 只需
17毫秒即可完成，这代表了超过六个数量级的加速。
话虽如此，问题本质上是指数级的，所以我们应该预期
\co{herd} 对更大的问题会表现出指数级的减慢。
这正是所发生的，例如，如果我们为每个写入 CPU 再添加四次写入，
如\cref{lst:formal:Expanded IRIW Litmus Test}所示，
\co{herd} 减慢了超过50,000倍，需要超过
15\emph{分钟}的 CPU 时间。
添加线程也会导致指数级的
减慢~\cite{PaulEMcKenney2014weakaxiom}。

尽管具有指数级的性质，PPCMEM 和 \co{herd} 都已被证明
对检查关键的并行算法非常有用，包括 x86 系统上的
队列锁交接。
\co{herd} 工具的弱点与 PPCMEM 类似，
后者在\cref{sec:formal:PPCMEM Discussion}中描述。
有一些隐蔽（但非常真实）的情况，PPCMEM 和
\co{herd} 工具不一致，截至2021年，其中许多但不是全部分歧
已被解决。

如果 litmus test 可以用 C 语言编写
（如\cref{lst:formal:Meaning of PPCMEM Litmus Test}）
而不是汇编语言
（如\cref{lst:formal:PPCMEM Litmus Test}），
那将会很有帮助。
这现在是可能的，将在以下各节中描述。

\subsection{Axiomatic Approaches 和锁}
\label{sec:formal:Axiomatic Approaches and Locking}

Axiomatic approaches 也可以应用于更高级别的语言
以及更高级别的同步原语，
如\cref{lst:formal:Locking Example}（\path{C-Lock1.litmus}）
中所示的基于锁的 litmus test 所例证的。
这个 litmus test 可以由
\IXalthmr{\acrf{lkmm}}{Linux kernel}{memory model}~%
\cite{Alglave:2018:FSC:3173162.3177156,LucMaranget2018lock.cat}建模。
如预期的那样，\co{herd} 工具的输出包含字符串 \co{Never}，
正确地表示 \co{P1()} 不能看到 \co{x} 具有值
1。\footnote{
	\co{herd} 工具的输出与 PPCMEM 兼容，
	所以可以查看
	\cref{lst:formal:PPCMEM Detects an Error,%
	lst:formal:PPCMEM on Repaired Litmus Test}
	中展示输出格式的示例。}

\begin{listing}
\input{CodeSamples/formal/herd/C-Lock1@whole.fcv}
\caption{锁示例}
\label{lst:formal:Locking Example}
\end{listing}

\QuickQuiz{
	你需要做什么才能在如
	\cref{lst:formal:Locking Example}所示的 litmus test 上运行 \co{herd}？
}\QuickQuizAnswer{
	获取 Linux 内核源代码的 v4.17（或更高）版本，
	然后按照 \path{tools/memory-model/README} 中的说明
	安装所需的工具。
	然后按照进一步的说明在你选择的 litmus test 上
	运行这些工具。
}\QuickQuizEnd

\begin{listing}
\input{CodeSamples/formal/herd/C-Lock2@whole.fcv}
\caption{有缺陷的锁示例}
\label{lst:formal:Broken Locking Example}
\end{listing}

当然，如果 \co{P0()} 和 \co{P1()} 使用不同的锁，如
\cref{lst:formal:Broken Locking Example}（\path{C-Lock2.litmus}）
所示，那么一切都不确定了。
在这种情况下，\co{herd} 工具的输出包含字符串
\co{Sometimes}，正确地表示使用不同的锁允许
\co{P1()} 看到 \co{x} 具有值1。

\QuickQuiz{
	为什么要费心直接建模锁？
	为什么不简单地用原子操作模拟锁？
}\QuickQuizAnswer{
	一个词，性能，如
	\cref{tab:formal:Locking: Modeling vs. Emulation Time (s)}中所示。
	第一列显示建模的 \co{herd} 进程数量。
	第二列显示直接在 \co{herd} 的 cat 语言中建模
	\co{spin_lock()} 和 \co{spin_unlock()} 时的 \co{herd} 运行时间。
	第三列显示使用 \co{cmpxchg_acquire()} 模拟 \co{spin_lock()}
	和使用 \co{smp_store_release()} 模拟 \co{spin_unlock()} 时的
	\co{herd} 运行时间，使用 \co{herd} 的 \co{filter} 子句
	拒绝未能获取锁的执行。
	第四列与第三列类似，但使用 \co{xchg_acquire()} 而不是
	\co{cmpxchg_acquire()}。
	第五列和第六列与第三列和第四列类似，
	但改用 \co{herd} 的 \co{exists} 子句来拒绝
	未能获取锁的执行。

\begin{table}
\rowcolors{10}{}{lightgray}
\small
\centering
\newcommand{\lockfml}[1]{\multicolumn{1}{c}{\begin{picture}(6,8)(0,0)\rotatebox{90}{#1}\end{picture}}}
\begin{tabular}{rrrrrr}
	\toprule
	& 建模 & \multicolumn{4}{c}{模拟} \\
	\cmidrule(l){2-2} \cmidrule(l){3-6}
	& & \multicolumn{2}{c}{\tco{filter}} & \multicolumn{2}{c}{\tco{exists}} \\
	\cmidrule(l){3-4} \cmidrule(l){5-6}
	\lockfml{进程数}
	&
	& \tco{cmpxchg}
	& \tco{xchg}
	& \tco{cmpxchg}
	& \tco{xchg}
	\\
	\cmidrule{1-1} \cmidrule(l){2-2} \cmidrule(l){3-4} \cmidrule(l){5-6}
	2 & 0.004 &  0.022 &  0.027 &   0.039 &   0.058 \\
	3 & 0.041 &  0.743 &  0.968 &   1.653 &   3.203 \\
	4 & 0.374 & 59.565 & 74.818 & 151.962 & 500.960 \\
	5 & 4.905 \\
	\bottomrule
\end{tabular}
\caption{锁：建模 vs.\@ 模拟时间（秒）}
\label{tab:formal:Locking: Modeling vs. Emulation Time (s)}
\end{table}

	还要注意，使用 \co{filter} 子句大约比使用
	\co{exists} 子句快两倍。
	这并不奇怪，因为 \co{filter} 子句允许
	早期放弃被排除的执行，而被排除的执行
	是锁被多个进程同时持有的执行。

	更重要的是，直接建模 \co{spin_lock()} 和 \co{spin_unlock()}
	比建模模拟的锁快5倍到超过两个数量级。
	这也不应该令人惊讶，因为直接建模提高了
	抽象级别，从而减少了 \co{herd} 必须建模的事件数量。
	因为 \co{herd} 所做的几乎所有事情都具有指数级的
	计算复杂度，事件数量的适度减少
	会产生指数级的运行时间缩减。

	因此，在 formal verification 中甚至比在并行编程本身中
	更要强调分而治之！！！
}\QuickQuizEnd

但锁不是唯一可以直接建模的同步原语：
下一节将介绍 RCU\@。

\subsection{Axiomatic Approaches 和 RCU}
\label{sec:formal:Axiomatic Approaches and RCU}

\begin{fcvref}[ln:formal:C-RCU-remove:whole]
Axiomatic approaches 也可以分析涉及
RCU 的 litmus test~\cite{Alglave:2018:FSC:3173162.3177156}。
为此，\cref{lst:formal:Canonical RCU Removal Litmus Test}
（\path{C-RCU-remove.litmus}）
展示了一个对应于经典的 RCU 介导的从链表中删除操作的
litmus test。
与锁 litmus test 一样，这个 RCU litmus test 可以
由 LKMM 建模，与建模 RCU 的模拟相比具有
类似的性能优势。
\Clnref{head}将 \co{x} 显示为列表头，初始
引用 \co{y}，后者在\clnref{tail:1}上被初始化为值 \co{2}。

\begin{listing}
\input{CodeSamples/formal/herd/C-RCU-remove@whole.fcv}
\caption{经典 RCU 删除 Litmus Test}
\label{lst:formal:Canonical RCU Removal Litmus Test}
\end{listing}

\Clnrefrange{P0start}{P0end}上的 \co{P0()}
通过将元素 \co{y} 替换为元素 \co{z}
（\clnref{assignnewtail}）来从列表中删除元素 \co{y}，
等待一个 \IX{grace period}（\clnref{sync}），
最后将 \co{y} 清零以模拟 \co{free()}（\clnref{free}）。
\Clnrefrange{P1start}{P1end}上的 \co{P1()}
在一个 RCU 读端临界区内执行
（\clnrefrange{rl}{rul}），
获取列表头（\clnref{rderef}），然后
加载下一个元素（\clnref{read}）。
下一个元素应该是非零的，即尚未被释放
（\clnref{exists_}）。
几个其他变量被输出用于调试目的
（\clnref{locations_}）。

\co{herd} 工具运行此 litmus test 时的输出包含
\co{Never}，表示 \co{P0()} 从不访问已释放的元素，
这是预期的。
同样如预期的那样，删除\clnref{sync}会导致 \co{P0()}
访问已释放的元素，如 \co{herd} 输出中的 \co{Sometimes}
所指出的。
\end{fcvref}

\begin{listing}
\ebresizeverb{.98}{\input{CodeSamples/formal/herd/C-RomanPenyaev-list-rcu-rr@whole.fcv}}
\caption{复杂的 RCU Litmus Test}
\label{lst:formal:Complex RCU Litmus Test}
\end{listing}

\begin{fcvref}[ln:formal:C-RomanPenyaev-list-rcu-rr:whole]
\ppl{Roman}{Penyaev}~\cite{RomanPenyaev2018rrRCU}提出的
一个更复杂的例子的 litmus test 如
\cref{lst:formal:Complex RCU Litmus Test}
（\path{C-RomanPenyaev-list-rcu-rr.litmus}）所示。
在这个例子中，读者（由\clnrefrange{P0start}{P0end}上的
\co{P0()} 建模）以轮询方式访问链表，
将指向最后访问的列表元素的指针
``泄漏''到变量 \co{c} 中。
更新者（由\clnrefrange{P1start}{P1end}上的 \co{P1()} 建模）
删除一个元素，注意不要干扰当前或未来的读者。

\QuickQuiz{
	等等！！！
	从 RCU 读端临界区泄漏指针
	不是一个严重的 bug 吗？？？
}\QuickQuizAnswer{
	是的，它通常是一个严重的 bug。
	然而，在这种情况下，更新者被巧妙地构造为
	正确处理此类指针泄漏。
	但请不要养成做这种事情的习惯，
	特别是在没有花大量心思使某种更常规的方法
	可行之前不要这样做。
}\QuickQuizEnd

\Clnrefrange{listtail}{listhead}定义了初始链表，
先从尾部开始。
在 Linux 内核中，这将是一个双向链接的循环列表，
但 \co{herd} 目前无法建模这样的结构。
因此策略是使用一个足够长的单链接线性列表，
使得永远不会到达末尾。
\Clnref{rrcache}定义了变量 \co{c}，用于在连续的
RCU 读端临界区之间缓存列表指针。

同样，\clnrefrange{P0start}{P0end}上的 \co{P0()} 建模读者。
这个进程建模了一对连续的读者以轮询方式遍历列表，
第一个读者在\clnrefrange{rl1}{rul1}上，
第二个读者在\clnrefrange{rl2}{rul2}上。
\Clnref{rdcache}获取缓存在 \co{c} 中的指针，
如果\clnref{rdckcache}看到指针为 \co{NULL}，
\clnref{rdinitcache}从列表开头重新开始。
无论哪种情况，\clnref{rdnext}前进到下一个列表元素，
\clnref{rdupdcache}将指向该元素的指针存储回
变量 \co{c} 中。
\Clnrefrange{rl2}{rul2}重复此过程，但使用
寄存器 \co{r3} 和 \co{r4} 而不是 \co{r1} 和 \co{r2}。
与\cref{lst:formal:Canonical RCU Removal Litmus Test}一样，
这个 litmus test 存储零来模拟 \co{free()}，所以
\clnref{exists_}检查这四个寄存器中是否有任何一个为
\co{NULL}，也就是零。

因为 \co{P0()} 从其第一个 RCU 读端临界区
向第二个泄漏了一个受 RCU 保护的指针，\co{P1()} 必须
非常小心地执行其更新（删除 \co{x}）。
\Clnref{updremove}通过将 \co{w} 链接到 \co{y} 来删除 \co{x}。
\Clnref{updsync1}等待读者，之后没有后续读者
可以通过链表找到通往 \co{x} 的路径。
\Clnref{updrdcache}获取 \co{c}，如果\clnref{updckcache}
确定 \co{c} 引用了新删除的 \co{x}，
\clnref{updinitcache}将 \co{c} 设置为 \co{NULL}，
\clnref{updsync2}再次等待读者，之后
后续读者无法从 \co{c} 获取 \co{x}。
无论哪种情况，\clnref{updfree}通过存储零到 \co{x}
来模拟 \co{free()}。

\QuickQuiz{
	\begin{fcvref}[ln:formal:C-RomanPenyaev-list-rcu-rr:whole]
	在\cref{lst:formal:Complex RCU Litmus Test}中，
	为什么读者不能在 \co{P1()} 在\clnref{updinitcache}
	上将 \co{c} 清零之前获取它，然后
	在它被清零之后将相同的值存储回 \co{c}，
	从而击败清零操作？
	\end{fcvref}
}\QuickQuizAnswer{
	\begin{fcvref}[ln:formal:C-RomanPenyaev-list-rcu-rr:whole]
	因为读者在\clnref{rdnext}上前进到下一个元素，
	从而避免存储指向与获取的元素相同的指针。
	\end{fcvref}
}\QuickQuizEnd

\co{herd} 工具运行此 litmus test 时的输出包含
\co{Never}，表示 \co{P0()} 从不访问已释放的元素，
这是预期的。
同样如预期的那样，删除任一 \co{synchronize_rcu()}
会导致 \co{P1()} 访问已释放的元素，如 \co{herd}
输出中的 \co{Sometimes} 所指出的。
\end{fcvref}

\QuickQuizSeries{%
\QuickQuizB{
	\begin{fcvref}[ln:formal:C-RomanPenyaev-list-rcu-rr:whole]
	在\cref{lst:formal:Complex RCU Litmus Test}中，
	为什么不在\clnref{updfree}之前只调用一次
	\co{synchronize_rcu()}？
	\end{fcvref}
}\QuickQuizAnswerB{
	\begin{fcvref}[ln:formal:C-RomanPenyaev-list-rcu-rr:whole]
	因为这会导致 \co{P0()} 访问已释放的元素。
	但不要只听我说，在 \co{herd} 中试一试！
	\end{fcvref}
}\QuickQuizEndB
%
\QuickQuizE{
	\begin{fcvref}[ln:formal:C-RomanPenyaev-list-rcu-rr:whole]
	同样在\cref{lst:formal:Complex RCU Litmus Test}中，
	\clnref{updfree}不能用 \co{WRITE_ONCE()} 而不是
	\co{smp_store_release()} 吗？
	\end{fcvref}
}\QuickQuizAnswerE{
	\begin{fcvref}[ln:formal:C-RomanPenyaev-list-rcu-rr:whole]
	这是一个很好的问题。
	截至2021年末，答案是``没人知道''。
	这在很大程度上取决于 \ARMv8 的条件移动指令的语义。
	在等待这些语义明确之前，\co{smp_store_release()}
	是安全的选择。
	\end{fcvref}
}\QuickQuizEndE
}

这些章节展示了 axiomatic approaches 如何能够成功地
在 C 语言 litmus test 中建模锁和 RCU 等同步原语。
从更长远来看，希望 axiomatic approaches 能建模
更高级别的软件工件，产生指数级的
验证加速。
这可能允许对更大的软件系统进行公理化验证，
也许结合分离逻辑中的空间同步
技术~\cite{AlexeyGotsman2013ESOPRCU,PeterWOHearn2001SeparationLogic}。
另一种选择是将布尔逻辑的公理投入使用，
如下一节所描述的。

% formal/sat.tex
% mainfile: ../perfbook.tex
% SPDX-License-Identifier: CC-BY-SA-3.0

\section{SAT 求解器}
\label{sec:formal:SAT Solvers}
%
\epigraph{靠启发式而活，因启发式而死。}{佚名}

任何具有有界循环和递归的有限程序都可以被转换为一个逻辑表达式，
该表达式可以用输入来表达程序的断言。
给定这样的逻辑表达式，如果能知道是否存在某种输入组合会触发其中一个断言，
那将非常有趣。
如果输入被表示为布尔变量的组合，
这就是所谓的 SAT，也称为可满足性问题。
SAT solvers 在硬件验证中被大量使用，这推动了巨大的进步。
20世纪90年代初的世界级 SAT solver 也许只能处理具有100个不同布尔变量的逻辑表达式，
但到21世纪10年代初，能处理百万变量的 SAT solver 已经
随处可得~\cite{Kroening:2008:DPA:1391237}。

\begin{figure}
\centering
\resizebox{2in}{!}{\includegraphics{formal/cbmc}}
\caption{\IXacr{cbmc} 处理流程}
\label{fig:formal:CBMC Processing Flow}
\end{figure}

此外，SAT solver 的前端程序可以自动将 C 代码翻译为逻辑表达式，
同时考虑断言并为数组越界等错误条件生成断言。
一个例子是 \IXacrl{cbmc}，即 \co{cbmc}，
它作为许多 Linux 发行版的一部分提供。
这个工具非常易于使用，只需 \co{cbmc test.c} 即可验证 \path{test.c}，
其处理流程如\cref{fig:formal:CBMC Processing Flow}所示。
这种易用性极其重要，因为它为将 formal verification 纳入回归测试框架
打开了大门。
相比之下，那些需要非平凡地翻译到专用语言的传统工具
只能用于设计时验证。

近年来，出现了能处理并行代码的 SAT solver。
这些求解器通过将输入代码转换为静态单赋值（SSA）形式，
然后生成所有允许的访问顺序来工作。
这种方法看起来很有前景，但它在实践中的效果如何还有待观察。
一个令人鼓舞的迹象是2016年将 \co{cbmc} 应用于 Linux 内核
RCU 的工作~\cite{LihaoLiang2016VerifyTreeRCU,Liang:2018:VTB,LanceRoy2017CBMC-SRCU}。
这项工作使用了 RCU 的最小配置，并使用少量线程验证了各种场景，
但仍然成功地读入了 Linux 内核的 C 代码并产生了有用的结果。
从 C 代码生成的逻辑表达式具有多达9000万个变量、4.5亿个子句，
占用数十 GB 的内存，
并且 SAT solver 需要多达80小时的 CPU 时间才能产生正确的结果。

然而，一个 Linux 内核开发者对于其代码已被自动验证的声明
可能有理由感到怀疑，
而且这样的怀疑者可以找到数十年来的许多
同道中人~\cite{DeMillo:1979:SPP:359104.359106}。
表达这种怀疑的一种有效方式是提供注入了 bug 的
被验证代码版本。
如果 formal verification 工具能找到所有注入的 bug，我们的开发者
可能会对该工具的能力增加一些信心。
当然，能找到开发者尚未察觉的有效 bug 的工具
可能会赢得更多信任。
这正是为什么有一个 \co{git} 归档包含20个分支的变异集，
每个分支可能包含一个注入到 Linux 内核 RCU 中的
bug~\cite{PaulEMcKenney2017VerificationChallenge6}。
任何拥有 formal verification 工具的人都被诚挚地邀请
在这组验证挑战上试用他们的工具。

目前，\co{cbmc} 能够找到一些注入的 bug，
但它尚未能找到 RCU 维护者
之前不知道的 bug。
尽管如此，仍有理由期望 SAT solver 有朝一日
能对寻找并行代码中的并发 bug 有所帮助。

% formal/stateless.tex
% mainfile: ../perfbook.tex
% SPDX-License-Identifier: CC-BY-SA-3.0

\section{无状态模型检查器}
\label{sec:formal:Stateless Model Checkers}
%
\epigraph{他列了一个清单，对它做了两次排列\dots}
	{改编自 Haven Gillespie 和 J. Fred Coots}

上一节描述的 SAT solver 方法非常方便且强大，但对所有可能执行
（包括状态）的完整跟踪可能会产生大量开销。
事实上，内存和 CPU 时间开销可能会严重限制可行验证的程序规模，
这引发了一个问题：不那么精确的方法是否可能在更大的程序中找到 bug。

\begin{figure}
\centering
\resizebox{2.1in}{!}{\includegraphics{formal/nidhugg}}
\caption{Nidhugg 处理流程}
\label{fig:formal:Nidhugg Processing Flow}
\end{figure}

虽然对这个问题的评判尚无定论，但 stateless model checker
（如 \IX{Nidhugg}~\cite{CarlLeonardsson2014Nidhugg}）在某些
情况下已经能处理更大的程序~\cite{SMC-TreeRCU}，并且具有
类似的易用性，如\cref{fig:formal:Nidhugg Processing Flow}所示。
此外，在某些 Linux 内核 RCU 验证场景中，Nidhugg 比
\co{cbmc} 快了一个数量级以上。
当然，Nidhugg 的速度和可扩展性优势与它不处理数据非确定性有关，
但这在这些特定的验证场景中不是一个因素。

尽管如此，与 \co{cbmc} 一样，Nidhugg 尚未能找到 Linux 内核 RCU
维护者之前不知道的 bug。
然而，它确实能够证明 Linux 内核 RCU 中的一个历史 bug
是由不同于维护者所认为的提交修复的，
这给了人们更多的希望，即像 Nidhugg 这样的 stateless model checker
有朝一日可能对在并行代码中找到并发 bug 有所帮助。


\section{总结}
\label{sec:formal:Summary}
%
\epigraph{西方思想一直聚焦于真-假；
	  现在是时候转向鲁棒-脆弱了。}
	 {Nassim Nicholas Taleb，概述}
% Full quote:
% Since Plato, Western thought and the theory of knowledge has focused on
% the notions of True-False; as commendable as that was, it is high time
% to shift the concern to Robust-Fragile, and social epistemology to the
% more serious problem of Sucker-Nonsucker.

本章描述的 formal verification 技术是验证小型并行算法的
非常强大的工具，但它们不应该是你工具箱中的唯一工具。
尽管数十年来一直关注 formal verification，测试仍然是大型并行软件
系统的验证主力~\cite{JonathanCorbet2006lockdep,DaveJones2011Trinity,PaulEMcKenney2016Formal}。

然而，情况并非永远如此是完全可能的。
要理解这一点，考虑到截至2017年，Linux 内核的实例估计超过
200亿个。
假设 Linux 内核有一个平均每百万年运行时间才会出现的 bug。
正如前一章末尾所指出的，这个 bug 将在安装基数上
每天出现超过50次。
但事实仍然是，大多数形式化验证技术只能用于
非常小的代码库。
那么，并发编程者该怎么办？

从找到第一个 bug、第一个相关 bug、最后一个相关 bug
和最后一个 bug 的角度来思考。

第一个 bug 通常通过检查或编译器诊断发现。
尽管日益复杂的编译器诊断构成了一种轻量级的 formal verification，
但人们通常不会这样看待它们。
这部分是由于一种奇特的从业者偏见，一方面说
``如果我在使用它，它就不可能是 formal verification''，
另一方面编译器诊断和验证研究之间存在很大的差距。

虽然第一个相关 bug 可能通过检查或编译器诊断找到，
但这两个步骤只找到拼写错误和误报的情况并不少见。
无论如何，大部分相关的 bug，即那些在生产中可能实际遇到的 bug，
通常会通过测试发现。

当测试由预期或实际的用例驱动时，
最后一个相关 bug 通过测试找到的情况并不少见。
这种情况可能会激发对 formal verification 的完全拒绝，
然而，不相关的 bug 有一个恼人的习惯，会在最不方便的时刻
突然变得相关，这要归功于黑帽攻击。
对于安全关键的软件，它似乎是总量中持续增长的一部分，
因此可能有强烈的动机来找到并修复最后一个 bug。
测试明显无法找到最后一个 bug，所以 formal verification 有
一个可能的角色，当然前提是 formal verification 能够
证明自己有能力承担这个角色。
正如本章所示，当前的 formal verification 系统
极为有限。

\QuickQuiz{
	但是，足够底层的软件难道不应该实际上
	对黑帽的利用免疫吗？
}\QuickQuizAnswer{
	不幸的是，不是。

	Paul E. McKenney 曾经认为 Linux 内核的 RCU
	对这种利用是免疫的，但 Row Hammer 的出现
	让他改变了看法。
	毕竟，如果黑帽能够攻击系统的 DRAM，
	他们就能攻击任何和所有底层软件，甚至包括 RCU\@。

	在2018年，这种可能性从理论推测的领域
	进入了客观现实的确凿领域~\cite{McKenney:2019:CRS:3319647.3325836}。
}\QuickQuizEnd

请注意，formal verification 通常比测试难用得多。
这在一定程度上是一个文化声明，有理由希望
随着熟悉度的增加，formal verification 会被认为更容易。
话虽如此，非常简单的测试框架就能在任意大的软件系统中
找到重大 bug。
相比之下，应用 formal verification 所需的工作量
似乎随着系统规模的增加而急剧增长。

尽管如此，我在近30年中偶尔使用 formal verification，
利用 formal verification 的优势，
即在设计时验证整体软件构造中小型复杂部分。
更大的整体软件构造当然是通过测试来验证的。

\QuickQuiz{
	鉴于 L4 微内核的完整验证，
	这种对 formal verification 的有限观点
	不是有点过时了吗？
}\QuickQuizAnswer{
	不幸的是，不是。

	L4 微内核的第一次完整验证是一项壮举，
	大量的博士生以非常低的每人速率
	手工验证代码。
	这种程度的工作量无法应用于大多数软件项目，
	因为变更速度太快了。
	此外，虽然 L4 微内核从 formal verification 的
	角度来看是一个大型软件工件，但与大量项目相比它很小，
	包括 LLVM、\GCC、Linux 内核、Hadoop、MongoDB
	以及许多其他项目。
	此外，这个验证确实有限制，研究人员
	坦然承认，值得赞扬：
	\url{https://docs.sel4.systems/projects/sel4/frequently-asked-questions.html\#does-sel4-have-zero-bugs}。

	虽然 formal verification 终于开始显示出一些
	前景，包括涉及更高自动化水平的更新的 L4 验证，
	但它目前在可预见的未来完全没有可能
	取代测试。
	虽然我非常希望在这一点上被证明是错误的，
	但请注意，这样的证明将以一个验证真实软件的
	真实工具的形式出现，而不是以大量
	激动人心的言辞出现。

	也许有一天 formal verification 将被大量用于
	验证，包括现在所知的回归测试。
	\Cref{sec:future:Formal Regression Testing?}审视了
	使这种可能性成为现实所需的条件。
}\QuickQuizEnd

最后一种方法是考虑
\cref{sec:debugging:Required Mindset}中的
以下两个定义及其隐含的结论：

\begin{description}[itemsep=0pt,labelindent=1em]
\item[定义：]	无 bug 的程序是平凡的程序。
\item[定义：]	可靠的程序没有已知的 bug。
\item[推论：]	任何非平凡的可靠程序至少包含
			一个尚未发现的 bug。
\end{description}

从这个角度来看，验证和确认方面的任何进步只能
有两种效果：
\begin{enumerate*}[(1)]
\item 平凡程序数量的增加，或
\item 可靠程序数量的减少。
\end{enumerate*}
当然，人类对多核系统和软件日益增长的依赖
为平凡程序数量的急剧增加提供了极大的动力。

然而，如果你的代码复杂到你发现自己过度依赖
formal verification 工具，你应该仔细重新考虑你的设计，
特别是如果你的 formal verification 工具要求你将代码
手动翻译到专用语言的话。
例如，\cref{sec:formal:Promela Parable: dynticks and Preemptible RCU}
中介绍的可抢占 RCU 的 dynticks 接口的复杂实现
结果证明有一个简单得多的替代实现，如
\cref{sec:formal:Simplicity Avoids Formal Verification}中所讨论的。
在其他条件相同的情况下，一个简单的实现远胜于
一个复杂实现的正确性证明。

对于那些从事 formal verification 技术和系统工作的人，
公开的挑战是证明这个总结是错误的！
为了协助完成这项任务，Verification Challenge~6 现已
发布~\cite{PaulEMcKenney2017VerificationChallenge6}。
来挑战吧！！！

\section{选择验证计划}
\label{sec:formal:Choosing a Validation Plan}
%
\epigraph{科学是楼上房间中的一流家具，
	  但前提是楼下要有常识。}
	 {Oliver Wendell Holmes，经修改}

你应该为你的项目使用什么样的验证？

正如软件领域特别是工程领域通常的情况一样，
答案是``视情况而定''。

请注意，运行测试或进行 formal verification
都不会改变你的项目。
充其量，这种努力通过定位一个随后被修复的 bug 而产生间接效果。
然而，修复一个 bug 可能会避免不便、经济损失、
财产损害，甚至生命损失。
显然，这种间接效果可能极其有价值。

不幸的是，正如我们所见，很难预测
给定的验证工作是否会发现重要的 bug。
因此很容易投入太少——甚至完全没有投入，
特别是如果开发估算过于乐观或预算意外紧张的话，
这些情况在现实世界的软件项目中几乎总会出现。

决定仍然投资于验证通常是由有强势个性的
经验丰富的人推动的。
但这并不能保证，因为其他利益相关者可能
同样具有强势的个性。
更糟糕的是，这些其他利益相关者可能会带来关于昂贵的
验证工作仍然让令人尴尬的 bug 逃逸到
最终用户的故事。
因此，虽然一个伤痕累累、满头白发、脾气暴躁的老手
可能会占上风，但也许一种更有组织的方法会更有用。

幸运的是，验证投资有一个严格的金融类比，
那就是保险单。

\IfTwoColumn{
\begin{figure*}
\centering
\resizebox{6in}{!}{\includegraphics{CodeSamples/formal/data/RCU-test-ratio.pdf}}
\caption{Linux 内核 RCU 测试代码}
\label{fig:formal:Linux-Kernel RCU Test Code}
\end{figure*}
}{}

保险单和验证工作都需要持续的前期投资，
两者都是对可能永远不会发生的灾难的防御。
此外，两者都有各种类型的免责条款。
例如，沿海地区的保险单可能不包括
海啸造成的损害，另一方面我们已经看到，
目前还没有任何验证方法能够找到
每一个 bug。

此外，在保险和验证方面都可能
过度投资。
仅举一例，一个消耗了整个开发预算的验证计划
将与一个覆盖太阳变成新星的保险单一样毫无意义。

一种方法是将软件预算的给定比例用于验证，
该比例取决于软件的关键性，
因此安全关键的航空电子软件可能将其预算的
更大比例用于验证，而家庭作业则不必如此。
在可用的情况下，应借鉴以往类似项目的经验。
然而，有必要构建项目，使验证投资在项目开始时就开始，
否则编码支出中不可避免的超支将挤占验证工作。

用经验丰富的人员来配置创业项目可能导致
验证工作的过度投资。
正如购买过多保险可能导致破产一样，
在测试上投入过多也可能扼杀一个项目。
这对于首创性项目尤其如此，因为尚不清楚
哪些用例将是重要的，在这种情况下，
对所有可能的用例进行测试将是对时间、
精力和资金的可能致命的浪费。

然而，随着创业项目支持的任务变得更加日常化，
用户通常对故障变得不那么宽容，从而增加了
对验证的需求。
管理这种投资转变可能极其具有挑战性，
特别是在用户不愿意或无法披露其用例确切性质的
太常见的情况下。
这时就需要从 bug 报告和与用户的讨论中
逆向工程用例变得至关重要。
随着这些用例被更好地理解，使用持续集成
可以帮助降低发现和修复所定位 bug 的成本。

\IfTwoColumn{}{
\begin{figure}
\centering
\IfEbookSize{
\resizebox{\onecolumntextwidth}{!}{\includegraphics{CodeSamples/formal/data/RCU-test-ratio-crop.pdf}}
}{
\resizebox{4.5in}{!}{\includegraphics{CodeSamples/formal/data/RCU-test-ratio-crop.pdf}}
}
\caption{Linux 内核 RCU 测试代码}
\label{fig:formal:Linux-Kernel RCU Test Code}
\end{figure}
}

\cref{fig:formal:Linux-Kernel RCU Test Code}展示了
一个软件项目验证使用的演变示例。
从该\lcnamecref{fig:formal:Linux-Kernel RCU Test Code}中可以看出，
Linux 内核的 RCU 直到 Linux 内核 v2.6.15 才有任何验证代码，
该版本的发布比2002年10月将 RCU 添加到 Linux 内核的
v2.5.43 版本晚了三年多。
测试套件在 Linux 内核 v2.6.19--v2.6.21 中达到了
其在总代码行数中的最高比例。
这个比例随着 v2.6.25 中接受用于实时应用的可抢占 RCU
而急剧下降。
这种下降是因为 RCU API 在可抢占和不可抢占的 RCU 变体中
是相同的。
这反过来意味着现有的测试套件适用于两种变体，
所以即使 Linux 内核的 RCU 代码显著扩展，
也不需要扩展测试。

\cref{fig:formal:Linux-Kernel RCU Test Code}中的后续条形图
显示 RCU 代码库显著扩展，但相应的
验证代码扩展得更为显著。
Linux 内核 v3.5 添加了对 \co{rcu_barrier()} API 的测试，
填补了测试覆盖中长期存在的空白。
Linux 内核 v3.14 添加了自动化测试和测试结果分析，
使 RCU 朝持续集成方向发展。
Linux 内核 v4.7 添加了 RCU 更新端原语的性能验证套件。
Linux 内核 v4.12 添加了 Tree SRCU，具有改进的更新端可扩展性，
v4.13 移除了旧的可扩展性较差的 SRCU 实现。
Linux 内核 v5.0 短暂地在 rcutorture 脚本目录中托管了
\path{nolibc} 库，然后它移到了 \path{tools/include/nolibc}
的长期归宿。
Linux 内核 v5.8 添加了 Tasks Trace 和 Rude 风格的 RCU\@。
Linux 内核 v5.9 添加了 \path{refscale.c} 读端性能测试套件。
Linux 内核 v5.12 和 v5.13 开始添加在运行时更改给定 CPU 的
回调卸载状态的能力，并添加了 \path{torture.sh} 验收测试脚本。
Linux 内核 v5.14 添加了分布式 rcutorture。
Linux 内核 v5.15 在 rcutorture 测试中添加了恶魔式 vCPU 放置，
这成功地定位了许多竞态条件。\footnote{
	技巧是将一对 vCPU 放在一个插槽上的同一核心内，
	同时将另一对放在某个其他插槽上的同一核心内。
	正如你可能从\cref{chp:Hardware and its Habits}中
	预期的那样，
	这会在不同的 vCPU 对之间产生不同的内存延迟
	（\url{https://paulmck.livejournal.com/62071.html}）。}
Linux 内核 v5.17 移除了 \co{RCU_FAST_NO_HZ} Kconfig 选项。
Linux 内核 v6.14 将 \co{kvfree_rcu()} 移到了 SLAB 中
（在 \path{mm/slab_common.c} 中）~\cite{20241212180208.274813-1-urezki}。
Linux 内核 v6.19 添加了一些 SRCU 变体，并在 \co{refscale}
测试套件中添加了若干原语。
许多其他更改可以在 Linux 内核的 \co{git} 归档中找到。
% rcutorture
% v2.6.15: First torture test
% v2.6.19: SRCU: Plugin architecture avoids test-code explosion.
% v2.6.19-21: Peak test fraction.
% v2.6.25: preemptible RCU, consistent API avoids added test code.
% v3.4: Add tests for RCU CPU stall warnings.
% v3.5: Add tests for rcu_barrier(). *
% v3.14: Add rcutorture scripting automating tests and results analysis. *
% v3.15: Add support for multiple torture-tests suites for locktorture.
% v3.16: Add support for conditional grace-period primitives.
% v4.7: Add update-side performance validation suite. *
% v4.12: Added Tree SRCU.
% v4.13: Removed non-Tree SRCU.
% v5.0: nolibc was briefly in the rcutorture scripting directory.
% v5.8: Added Tasks Trace RCU and Rude RCU.
% v5.9: Added refscale.c.
% v5.12: Added torture.sh.  Added runtime adjustment of rcu_nocbs CPUs.
% v5.13: More rcu_nocbs CPUs, polls grace-period primitives, remote rcutorture.
% v5.14: Added demonic affinity to rcutorture scripting.
% v5.17: Added RCU callback (de)offloading optimizations and per-CPU
%	 tracking of RCU Tasks callbacks.
% v5.18: Improved RCU priority boosting and \co{rcu_barrier()} no longer
%	 blocks CPU-hotplug operations.
% v5.19: Enabled milliseconds-scale real-time response from
%	 \co{synchronize_rcu_expedited()} and also dynamically allocated
%	 and sized Tree SRCU's \co{srcu_node} array.  This last saves
%	 significant memory on kernels sized for thousands of CPUs
%	 that typically run on much smaller systems.

我们已经确定验证预算因项目而异，
也会在任何给定项目的生命周期中变化。
但是验证投资应该如何在测试和
formal verification 之间分配呢？

这个问题正在自然地被回答，因为编译器采用了日益
激进的 formal verification 技术到其诊断中，
而且 formal verification 工具继续成熟。
此外，Linux 内核的 lockdep 和 KCSAN 工具说明了
将 formal verification 技术与运行时分析相结合的优势，
如\cref{sec:debugging:Assertions}中所讨论的。
其他组合技术分析从执行中收集的
轨迹~\cite{DanielBristot2019RTtrace}。
目前，最佳实践是首先关注测试，并将
formal verification 的显式工作保留给项目中
测试不能很好服务的那些部分，以及对鲁棒性有
特殊需求的部分。
例如，Linux 内核的 RCU 主要依赖测试，但
如本章所讨论的，偶尔使用 formal verification。

总之，为并发软件选择验证计划仍然更多的是一门
艺术而非科学，更不用说工程领域了。
然而，有充分的理由期望日益严格的方法
将继续变得更加普遍。

\QuickQuizAnswersChp{qqzformal}
